\pdfoutput=1
\documentclass[11pt]{amsart}
\usepackage{amsmath,amssymb,amsthm,mathtools}
\usepackage[margin=1.2in]{geometry}
\usepackage{microtype,xcolor}
\usepackage[colorlinks=true,linkcolor=blue,citecolor=blue,urlcolor=blue]{hyperref}
\hypersetup{pdfauthor={Bernd Johannes Wuebben},pdftitle={Spin structures and local Gromov--Witten invariants of curves}}
\setlength{\emergencystretch}{2em}
\newtheorem{theorem}{Theorem}[section]
\newtheorem{proposition}[theorem]{Proposition}
\newtheorem{lemma}[theorem]{Lemma}
\newtheorem{corollary}[theorem]{Corollary}
\theoremstyle{definition}
\newtheorem{definition}[theorem]{Definition}
\newtheorem{example}[theorem]{Example}
\theoremstyle{remark}
\newtheorem{remark}[theorem]{Remark}
\newcommand{\PP}{\mathbb P}
\newcommand{\CC}{\mathbb C}
\newcommand{\ZZ}{\mathbb Z}
\newcommand{\QQ}{\mathbb Q}
\newcommand{\RR}{\mathbb R}
\newcommand{\cO}{\mathcal O}
\DeclareMathOperator{\sgn}{sgn}
\raggedbottom
\setcounter{tocdepth}{1}
\title[Spin structures and local curve invariants]{Spin structures and local\\ Gromov--Witten invariants of curves}
\author{Bernd Johannes Wuebben}
\subjclass[2020]{53D45 (primary); 14H30, 14H60, 14N35 (secondary)}
\keywords{Local Gromov--Witten invariants, spin structures, relative Gromov--Witten invariants, branched covers}
\date{November 24, 2026\\This version: November 24, 2026}
\begin{document}
\begin{abstract}
We compute the local Gromov--Witten invariants of a curve of genus at least two with positive maximal flat symplectic normal holonomy, in every covering degree and source genus. The orthogonal bundle of the harmonic reduction determines which partitions occur. When it is orientable, the series agrees with the antidiagonal local-curve series of Bryan and Pandharipande. When it is nonorientable, only self-conjugate partitions contribute. The determinant sign on the embedded curve multiplies the degree variable in this sum. The restriction on partitions comes from the change of orientation at the relative contacts when two spin structures are glued across a nonseparating node. On the closure of the locus of triple covers with four simple branch points over a genus-two curve, we compute the intersection matrix of the ten classes measuring the difference between the real spin Euler class and the corresponding complex Chern class, one for each even spin structure.
\end{abstract}
\maketitle
%\newpage
%\tableofcontents
%\newpage
\section{The local contribution of a curve}\label{sec:partition}

How much does the normal operator of an embedded curve tell us about
the contribution of its multiple covers? Even if every cover is
isolated in the normal directions, counting the unramified covers with
their determinant signs gives only the first possible source genus in
each degree. Higher genera allow branching, nodes and contracted
components. We compute their combined contribution for a curve with
positive maximal flat symplectic normal holonomy. Already in degree
two, changing one of two identical spin structures can make the entire
disconnected contribution vanish. The reason appears when the target
is cut at a nonseparating node. Reversing one spin gluing tensors the
boundary character with the sign representation. It exchanges a Young
diagram with its transpose, so only self-conjugate diagrams contribute
when the two ends are glued back together.

Let \(C\) be a smooth projective complex curve of genus \(g\ge2\), and
let \(\rho:\pi_1(C)\to{\rm Sp}(4,\RR)\) be the holonomy of a flat
symplectic four-plane bundle. We use the convention that \(\rho\) is
\emph{positive maximal} if its compatible complex rank-two bundle
has degree \(2g-2\). In particular the total space has first Chern
number zero on its zero section.
The harmonic reduction gives a holomorphic bundle \(E\) and a
nondegenerate symmetric pairing \(\gamma:E\otimes E\to K_C\).
Choose a theta characteristic \(a\), that is, a holomorphic square
root of \(K_C\). Untwisting the pairing gives the complex orthogonal
bundle \(F=E\otimes a^{-1}\), the underlying bundle of the Cayley
correspondence of Bradlow, Garc\'ia-Prada and Gothen \cite{bgg}. Its first Stiefel--Whitney class \(w_1(F)\) measures
whether its transition functions can be chosen in \({\rm SO}(2,\CC)\).
In that case its two isotropic lines give \(F=M\oplus M^{-1}\);
write \(s=|\deg M|\). For \(w_1(F)\ne0\) monodromy interchanges the
isotropic directions, so an antidiagonal action scaling the two
lines separately need not exist.

A Hermitian metric turns the symmetric pairing into an antilinear map
\(R_\gamma:E\to\Lambda^{0,1}C\otimes E\) by taking the metric dual and
conjugating the form factor. The positive normal operator is
\(D=\bar\partial_E+R_\gamma\). In a split model \(E=a_1\oplus a_2\)
with \(a_i^2=K_C\), the pairing is the sum of the two square maps and a
product metric gives two independent real spin operators. For
\(w_1(F)\ne0\), every unmarked deformation type has a representative
of this form with \(a_1\ne a_2\). Conversely, the direct sum of the two uniformizing spin lifts realizes each such pair. Write \(\epsilon\) for the
determinant sign of \(D\) on \(C\), with orientation continued from the
complex Dolbeault operator. On such a representative
\(\epsilon=(-1)^{h^0(a_1)+h^0(a_2)}\).
Isolation, reduction to these representatives and invariance under
normal rescaling are proved in
\cite[Lemma~14.1, Propositions~14.2 and~15.4, and Appendices~D--E]{WueUnramified}.

Let \(\overline{\mathcal M}_k(C,m)\) be the proper stack of connected
unmarked stable maps of arithmetic genus \(k\) and degree \(m>0\) to
this fixed \(C\). Over a map \(f\), the antilinear term of \(D_f\)
uses \(df\), and sections on the branches of a node must agree.
The kernel is zero on every nonconstant component by the spin energy
identity. Matching then kills the constants on every contracted
subcurve. Its oriented real cokernel \(V_D\) consequently has rank
\(4(k-1-m(g-1))\), the real virtual dimension of the stack.
The ordinary connected local invariant and its disconnected series are
\begin{equation}\label{eq:fixed-target}
 L_{m,k}=\int_{[\overline{\mathcal M}_k(C,m)]^{\rm vir}}e(V_D),
 \qquad
 Z_D(Q,v)=\exp\!\left(\sum_{m\ge1,\,k}L_{m,k}Q^m v^{2k-2}\right).
\end{equation}
The sum starts at \(k=1+m(g-1)\) by Riemann--Hurwitz on the
normalization. Every connected component counted by \(Z_D\) has
positive degree. At the least genus each connected cover contributes its determinant sign divided by its deck order. Giving each cover the positive reciprocal of its deck order instead would recover the ordinary unramified Hurwitz number.

Bryan and Pandharipande computed the residue invariant of a direct sum
of two holomorphic line bundles with opposite torus weights
\cite[Corollaries~7.2--7.3]{bp08}. Their answer is a sum over all
partitions of the covering degree. Our first formula identifies the
real normal Euler integral with their residue invariant in the
orientable case. For a nonorientable \(F\), the sewing of the normal
operators changes that sum.

For a box \(u\) in a Young diagram \(\lambda\), write \(h(u)\) for
one plus the number of boxes strictly to its right or below it in its
row or column. Put \(a(\lambda)=\sum_{(i,j)\in\lambda}(j-i)\), the
total content. With \(t=e^{iv}\), let \(X_j=2-t^j-t^{-j}\) and
\(b_0=g-1\). Thus \(X_j=4\sin^2(jv/2)\).

\begin{theorem}\label{thm:spin}
For the positive maximal normal operator above, in every degree and
source genus the disconnected fixed-target series is
\[
 Z_D(Q,v)=
 \begin{cases}
 \displaystyle\sum_{\lambda}(-1)^{s|\lambda|}Q^{|\lambda|}
 t^{s a(\lambda)}\prod_{u\in\lambda}X_{h(u)}^{b_0},
       &w_1(F)=0,\\[5pt]
 \displaystyle\sum_{\lambda=\lambda^\top}(\epsilon Q)^{|\lambda|}
 \prod_{u\in\lambda}X_{h(u)}^{b_0},&w_1(F)\ne0.
 \end{cases}
\]
The empty diagram contributes one. The same series gives the local
invariants of a marked horizontal section of a flat symplectic
four-plane germ with this derivative holonomy and the split horizontal
and compatible vertical almost-complex structures.
\end{theorem}

The obstruction to the second comparison occurs at the seam of a
single handle. Pinch a nonseparating circle on \(C\), and choose the
two spin structures to agree on the normalization and to differ in the
identification of the two end fibers. For a relative map, the contact
orders make a partition \(\mu\) of \(m\). The mixed spin Euler integral
acquires the sign \(\epsilon^m(-1)^{m-\ell(\mu)}\) relative to the
square for a repeated spin. The second factor is the sign of a
permutation of cycle type \(\mu\). In the character basis, tensoring
by sign interchanges a partition and its transpose. Gluing the handle
takes a trace. A partition contributes to that trace exactly when
it is fixed by transposition. The orientation calculation has to
include even contacts and contracted components to give this sign
on the full relative virtual class.

Self-conjugate diagrams also occur in the local Real Gromov--Witten
theory of Georgieva and Ionel \cite[Theorem~1.1 and Corollary~9.14]{gi}.
Their domain and target carry antiholomorphic involutions and the
maps intertwine them. For their canonically oriented Real
Calabi--Yau series \(R_G(q,u)\) on a connected symmetric curve of
genus \(G\), the formal expression in the second line is precisely
\(R_{2g-1}((-1)^{g-1}\epsilon Q,iv)\). Indeed the crosscap sign in
their formula is raised to the even power \(2g-2\), and each
hyperbolic sine acquires a factor of \(i\). Our moduli in
\eqref{eq:fixed-target} are ordinary complex stable maps; the
orientation change comes from the \emph{normal spin matching} over
one node. In particular, taking the logarithm here selects connected
ordinary domains. In the Real theory the logarithm also counts pairs
of components exchanged by the involution \cite[(2.16)]{gi}.

For a concrete consequence, put \(x=X_1\). A balanced genus-\(j\)
factor is the disconnected antidiagonal local series of a curve of
genus \(j\) with both normal lines of degree \(j-1\); denote it
by \(B_j(Q,v)\). Its logarithm starts with \(Qx^{j-1}\).
The \emph{elementary multiplicities} \(n_{d,j}\) of a compactly
realized isolated germ are the integers in the factorization
\(Z_D=\prod_{d,j} B_j(Q^d,v)^{n_{d,j}}\), as justified by
\cite[Proposition~3.1]{WueUnramified}. The balanced local factor is the elementary curve of Ionel and Parker under the comparison of \cite[\S3]{WueProduct}. These multiplicities record a genus of the elementary
curve and differ from the integrals at a specified source genus.

\begin{corollary}\label{cor:cubic}
For two distinct spin structures of the same parity on a genus-two
target, the connected degree-three series is
\(28x^3/3-6x^4+x^5\). Its nonzero degree-three elementary
multiplicities in increasing genus are
\[
 (n_{3,4},n_{3,5},n_{3,6},n_{3,7})=(-64,64,-20,2).
\]
\end{corollary}
\begin{proof}
There is no self-conjugate partition of two. In degree three there
is just \((2,1)\), with hooks \(3,1,1\). Thus the logarithm gives
\(X_3x^2+x^3/3\), with \(X_3=9x-6x^2+x^3\). The only proper
divisor of three is one. Subtracting the degree-three term of the
balanced genus-two logarithm,
\(220x^3/3-70x^4+21x^5-2x^6\), gives the four stated integers.
\end{proof}

The degree-three answer is the same for any two distinct even
spin structures on a given genus-two curve. It gives the integral
of a product, summed over all stable maps. To recover the individual
classes one must locate them on the moduli space. We identify the
cokernel of a single spin operator with a real form of the residue
space at the ramification points. This gives a signed divisor formula
on the surface of covers with two simple branch points. On the
four-dimensional closure of the cubic covers with four simple branch
points the individual spin corrections have square \(6\) and mixed
pairing \(-1/3\) for any two of the ten distinct even spins
(Theorem~\ref{thm:gram}). Their common support is detected by the
difference of two squared Szeg\H{o} kernels. These intersections
explain what the partition sum forgets about an individual square
root of \(K_C\).

\section{Relative degeneration and the sign of a handle}\label{app:spinproof}

We prove Theorem~\ref{thm:spin}. Deformation of the metric and the
normal rescaling theorem in \cite[Proposition~15.4 and Appendices~D--E]{WueUnramified} identify the
fixed-target integral for any of the stated germs with its value at a
harmonic linear representative. At each fixed degree and source genus
this continuation takes place over the compact stable-map stack. The
vertical area identity on positive normalizations, followed by
matching on contracted components, keeps the normal family injective.

\subsection{The orientable normal bundle}
In the splitting \(E=L_+\oplus L_-\), the two degrees are
\(b_0+s\) and \(b_0-s\). Choose an orthogonal product metric.
Let \(\pi\) and \(f\) denote the universal source and universal map
over \(\overline{\mathcal M}_k(C,m)\), and put \(W_\pm=-R\pi_*f^*L_\pm\),
\(n=k-1-mb_0\). The antilinear term of \(D\) exchanges the lines;
the family from Dolbeault to the positive operator is complex-linear
in the complex structure \((i,-i)\). At its positive end the kernel is
zero. At its Dolbeault end the orientation change on the second
summand is \((-1)^{n+ms}\), since this is the parity of its
Dolbeault index. Stabilization by forms supported away from the nodes
continues the family index, including the complex equations at the
nodes, across jumps in the individual cohomology bundles. Thus
\begin{equation}\label{eq:oriented}
 e(V_D)=(-1)^{n+ms}c_{2n}(W_++\overline W_-).
\end{equation}
The right side is the cohomological degree \(2n\) term (in complex
degrees) of the equivariant Euler product for weights \(+1,-1\).
All other terms integrate to zero. Bryan--Pandharipande's
antidiagonal formula \cite[Corollary~7.2]{bp08} now gives the first
line of the theorem. Exchanging the weight convention transposes all
diagrams and preserves the sum.

\subsection{Two smoothings of the spin pair}
Pinch a nonseparating circle on \(C\); denote the normalization of
the central target by \(B\) and its ordered inverse images of the
node by \(p,q\). For each \(\epsilon\) we need one smoothing of a
pair of spins whose ratio has nontrivial holonomy around the loop
of the dual graph. The unmarked component classification in
\cite[Appendix~E]{WueUnramified} then applies to every pair of that sign.
For \(\epsilon=1\) use the exceptional spin at an index-two node.
The two spin structures restrict to the same theta characteristic
on \(B\) and differ by a sign in sewing at \(q\). Over an exceptional
rational component the spin has degree one, so its sections interpolate
the two end values uniquely. Both extensions thus have the same
parity. For \(\epsilon=-1\), take the two locally free extensions of
\(A^2=K_B(p+q)\) at an ordinary node. Riemann--Roch gives
\(h^0(A)-h^0(A(-p-q))=1\); the line of end values is isotropic for
the sum of squared residues, by the residue theorem. Hence exactly one
of the two isotropic clutch lines admits the extra section. The two
parities differ. The roots smooth after a local base change, by
smoothness of the relative Picard space at a line and étaleness of
squaring; parity is preserved in spin smoothing \cite[\S6]{cornalba}. In both cases their
ratio has trivial inertia and its central double torsor is connected
by the nontrivial clutch. This persists under smoothing. See
Calculation~D1 of \cite{WueSupplement} for the square frames.
Use \(d=2\) for the exceptional spin and \(d=1\) for the locally free
spin.

After \(\delta=z^2\), resolve the local smoothing \(UV=\delta\)
before taking the stabilizer quotient. The two central components
are the \(d\)-th root of \(B\) at \(p,q\) and the two-ended
rational component \([\PP^1/\mu_d]\) with its rotation action.
Both have coarse covering degree \(m\), by intersection with the
exceptional component in the surface. We use the simultaneous
relative expansions at \(p,q\) of \cite[\S\S4.4--4.7]{af}.
In a frame whose square is the logarithmic differential, take the
antilinear term proportional to \(U\,d\bar U-V\,d\bar V\) on the neck.
Its coefficient relative to the conjugate logarithmic differential is
\(|U|^2+|V|^2\) on a smooth fiber and can be joined by positive
interpolation to the metric off the neck. Pullback by contraction
makes it zero on the inserted components.

Stokes' identity on a main normalization gives nonnegative interior
energy and opposite end residues weighted by the coarse contact
orders; orbifold integration cancels the uniformizer index. The
square of a holomorphic solution on an inserted component is
constant. Conservation of contact order thus cancels the end
residues on the glued source. The solution vanishes on each
nonconstant main component, then on all components by matching.
In particular contracted subcurves of arbitrarily high genus do not
create a kernel. Stabilization with forms off the nodes gives an
oriented cokernel bundle on the plumbing charts. The antilinear term
is the pullback of a bounded ambient coefficient on each half-neck,
so the nodal estimates of \cite[Appendix~B]{bai} apply with this term.
The family Euler product restricts under virtual Gysin to the paired
Euler insertion on the central splitting. The expanded family and
virtual restriction are described in
Calculations~D6--D7 of \cite{WueSupplement}.

\subsection{Eliminating the two-ended component}
On a connected exceptional relative source of arithmetic genus \(j\)
let \(S\) denote the pulled-back sign line. Its space of constants
is a line or zero according as its double torsor is trivial or
connected; for \(d=1\) it is always a line. These alternatives are
locally constant on the relative stack. The normalization sequence
gives the two real cokernels a common complex quotient
\(A=R^1\pi_*S\). Mumford's Hodge relation on the coarse double torsor,
divided by the same relation on its invariant summand, gives
\(c(A)c(A^\vee)=1\). In positive rank \(c_{\rm top}(A)^2=0\), so the
paired insertion vanishes. In rank zero eliminate the surjective
exceptional Dolbeault equations. The remaining equations involve the
main operator, its evaluations at the ends and one constant for each
trivial-torsor vertex. Pass to the finite étale cover framing the
invariant end lines by square-compatible frames. The residual choices
are finite signs, and on every such sector the insertion is pulled
back from the main relative space. Finite transfer thus reduces the
exceptional integral to its virtual class paired against \(1\).

The log tangent of the two-ended target is trivial. A connected
positive-degree exceptional source with \(a\) ends meets both colors,
so its virtual dimension \(2j-2+a\) is nonnegative. Integration against
\(1\) vanishes in positive dimension. Only genus zero with two ends
survives. Stability and Riemann--Hurwitz then leave the power map
of degree \(k\), branched only at the two ends; unbranched rubber
levels give the same relative map. The two profiles agree, say
\(\mu=(k_1,\ldots,k_r)\). The following multiplicity allows repeated
parts as well as even contacts.
\begin{lemma}\label{lem:twoends}
Suppose the surviving maps over the two-ended component are power
cylinders with coarse contact profile \(\mu=(k_i)\) and the insertion
is pulled back from the main relative factor. For one prescribed
matching of the separately labelled ends the weight is \(\prod_i k_i\).
Against the main integral with both lists unlabelled it is
\(z_\mu=|\operatorname{Aut}\mu|\prod_i k_i\). This holds for target
index \(d=1\) and \(d=2\).
\end{lemma}
\begin{proof}
At one flag of coarse order \(k\) put \(a=(d,k)\), \(f=d/a\)
and \(c=k/a\). The primitive representable source index is \(f\);
in a uniformizer the contact has order \(c\). The inverse-band
matching diagonal of the rigidified evaluation sector
\(B(\mu_d/\mu_f)\) has degree \(a\): its fiber is the torsor of
isomorphisms between two given evaluations. Choose one auxiliary
twist \(R\) divisible by every \(c\). The target index becomes
\(dR\), the source index \(dR/k\), and the map of rigidified
sectors to the primitive sector is a \(\mu_c\)-gerbe. Comparison
of the diagonals contributes \(c\) per flag. The single target
gluing gerbe of degree \(1/R\) and the central multiplicity \(R\)
cancel on the common expansion base. Thus the flag factor is
\(ac=k\) \cite[\S\S5.5--5.10, Lemma~5.15 and Proposition~5.18]{af}.

With \(N=dR\), the cylinder is the index-\(k\) rotation quotient
\([{\PP}^1/\mu_{N/k}]\to[{\PP}^1/\mu_N]\). The simply connected
atlas gives a unique connected cover for this subgroup. Its log
differential is an isomorphism, and its deck group has order \(k\).
The two matching factors and its virtual degree give \(k^2/k=k\).
In particular for \(d=2\) and \(k\) even the two relative spin
clutches each carry weight \(k/2\): the deck action changes both
primitive signs together, and the two resulting orbits have
stabilizer of order \(k/2\). The pullback assumption gives the
same insertion on these two orbits.

Labelling equal cylinders introduces the divisor
\(A=|\operatorname{Aut}\mu|\) once, simultaneously on both ends.
The unlabelled main integral divides the separately labelled integral
by \(A^2\). Their ratio gives \(A\prod_i k_i\), as asserted.
\end{proof}

\subsection{The orientation at a contact}
In the primitive chart the sign line has an invariant matching fiber
at order \(k_i\) exactly when \(d\mid k_i\). Write \(I=\{i:d\mid k_i\}\).
On a surviving cylinder in \(I\), elimination of its constant at the
first end leaves the complex equation \(q_i-\sigma_i p_i=0\), where
\(p_i,q_i\) are the end values in square frames and \(\sigma_i=\pm1\).
Changing the target clutch negates each \(\sigma_i\).
On the cut main source Stokes' boundary inequality, after a common
phase rotation of the frames, is
\[
 \sum_{i\in I} k_i\operatorname{Re}(p_i^2-q_i^2)\ge0
\]
with equality only for the zero solution of the main equations.
For \(0\le u\le1\) impose instead
\[
 \operatorname{Re}p_i-u\sigma_i\operatorname{Re}q_i=0,\qquad
 \operatorname{Im}q_i-u\sigma_i\operatorname{Im}p_i=0 .
\]
Each term in the boundary sum then equals
\(k_i(u^2-1)((\operatorname{Re}q_i)^2+(\operatorname{Im}p_i)^2)\le0\).
The augmented operators remain injective throughout. Their endpoint
at \(u=0\) is independent of the clutch. At \(u=1\) the transformation
from the original complex row to the new two real rows has
determinant \(-\sigma_i\). Thus each changed invariant matching
fiber reverses the orientation once. For \(d=2\) the number of even
parts has parity \(m-r\); for \(d=1\) all \(r\) parts contribute.
The mixed insertion equals the repeated square multiplied by
\begin{equation}\label{eq:clutch}
 (-1)^{|I|}=\epsilon^m(-1)^{m-r}.
\end{equation}
Before this comparison one must keep the two spin matchings at an
even contact separate. After it the sign depends only on \(\mu\),
so finite transfer preserves the formula.

\subsection{The repeated square and the trace}
Let \(H\) denote the spin pulled back to the closed source after the
cylinders have been attached, and \(W_H=-R\pi_*H\). The square of
the Euler class of an oriented real \(2n\)-plane bundle is its top
Pontryagin class. Continuation to Dolbeault gives
\[
 e(V_H)^2=(-1)^n c_{2n}(W_H+\overline W_H).
\]
We compare this square with an ordinary index on the coarse source.
Choose a locally free spin \(L\) on the ordinary nodal target. Pull
it back to the twisted target. Its square agrees with that of \(H\)
in logarithmic frames, so \(F'=HL^{-1}\) on the target is a sign
line. There \([H]-[L]=[F']-1\) in the perfect \(K\)-group: express
\(L\) by divisors off the node and tensor the resulting point
sheaves by \(F'\). Pull this identity back along the universal map
before pushing it down to the parameter space. The zeroth
cohomologies of the sign and trivial lines have finite monodromy;
the positive even Chern characters of their first cohomologies
vanish by the Hodge relation on the coarse double torsor and its
invariant summand. Conjugation cancels the odd Chern characters.
Consequently
\(c_{2n}(W_H+\overline W_H)=c_{2n}(W_L+\overline W_L)\) rationally,
including on the twisted sectors. On normalizing the \(r\) power
cylinders the two complex rows at each cylinder cancel its single
constant, leaving \(W_L=W_B+\mathcal O^r\) after framing. Here \(W_B\)
is the ordinary relative index on the main coarse source. It
descends with the relative virtual class by
\cite[Theorem~2.1]{acw} and \cite[Theorem~1.1.2]{amw}.

For equivariant weight \(s_0\), the inverse antidiagonal pairing on
a profile of length \(r\) is \(z_\mu(-s_0^2)^r\);
the required degree of the equivariant Euler product of \(W_B\)
contributes \((-1)^{n-r}s_0^{-2r}c_{2n}(W_B+\overline W_B)\).
Their product is \(z_\mu\) times the repeated square. The factor
\(z_\mu\) agrees with Lemma~\ref{lem:twoends} on unlabelled relative
maps. This establishes the normalization and the orientation for the
full contraction, including both clutch types. The arithmetic genus
of the disconnected cut source, defined by \(1-\chi(\mathcal O)\),
is \(k-r\). The relative shift contributes \(2r\) to its genus
exponent, restoring \(2k-2\).

Let \(A_m\) be the balanced two-point operator of genus \(g-1\) and
levels \((b_0,b_0)\) in the antidiagonal relative theory. The
degree-\(m\) coefficient of \(Z_D\) is now
\(\epsilon^m\operatorname{Tr}(\Omega A_m)\), where
\(\Omega e_\mu=(-1)^{m-\ell(\mu)}e_\mu\) in the relative basis
indexed by contact profiles. In the idempotent indexed by
\(\lambda\), the coefficient of \(e_\mu\) is proportional to
\(\dim(\lambda)\chi^\lambda(\mu)\), with the remaining factor
depending only on \(m,\mu,s_0\). Tensoring the character with sign
therefore gives \(\Omega v_\lambda=v_{\lambda^\top}\), with unit
coefficient. The eigenvalue of \(A_m\) is
\(\prod_{u\in\lambda}X_{h(u)}^{b_0}\) by the handle and level
eigenvalues in \cite[\S7]{bp08}. Taking the trace proves the second
line of Theorem~\ref{thm:spin}.

\section{The largest genus in the elementary factorization}\label{sec:highest}

The source-genus expansion of a fixed positive-degree term of
\eqref{eq:fixed-target} generally has infinitely many coefficients.
Its elementary multiplicities have finite support. The last nonzero
genus can be read from the two extreme diagrams in the antidiagonal
formula. Write \(n_{m,j}\) for the multiplicities defined after
Theorem~\ref{thm:spin}, using the compact realization and fixed-target
comparison of \cite[Proposition~3.1 and its proof]{WueUnramified}.

\begin{theorem}\label{thm:highest}
Suppose the hypotheses of Theorem~\ref{thm:spin} hold.
If \(w_1(F)=0\) and \(s>0\), the largest \(j\) with
\(n_{m,j}\ne0\) is \(1+D_m\) in each degree \(m\ge1\), where
\[
 D_m=mb_0+(b_0+s)\binom m2,\qquad
 n_{m,1+D_m}=(-1)^{m(b_0+s)+D_m}.
\]
For \(w_1(F)=0\) and \(s=0\), the only elementary multiplicity is
\(n_{1,g}=1\). For \(w_1(F)\ne0\), one has \(n_{1,g}=\epsilon\);
for \(m\ge2\) the largest genus is \(1+D_m\), with
\[
 D_m=b_0\binom{m+1}2,\qquad
 n_{m,1+D_m}=-2\epsilon(-1)^{mb_0+D_m}.
\]
\end{theorem}
\begin{proof}
Use Laurent degree in \(t\), so that \(X_i\) has degree \(i\).
The sum of the hooks of \(\lambda\vdash m\) is \(m\) plus the number
of pairs of boxes in a common row or column. It is at most
\(\binom{m+1}2\), with equality only for a single row or column.
The total content is at most \(\binom m2\), attained only by the
row. For \(s>0\) the row alone therefore contributes the top term
of degree \(D_m\), with coefficient \((-1)^{m(b_0+s)}\).
When \(s=0\), Theorem~\ref{thm:spin} gives \(Z_D=B_g\).

For \(w_1(F)\ne0\) and \(m>1\) neither extreme diagram is
self-conjugate. Thus the numerator in degree \(m\), and hence the
degree-\(m\) term in \(\log Z_D\), has Laurent degree strictly less
than \(D_m\): in passing to the logarithm, products from smaller
degrees are strictly below the bound by strict superadditivity.
For the same reason the top term in the balanced genus-\(g\) logarithm
is contributed by the row and column alone, with coefficient
\(2(-1)^{mb_0}\). It enters the elementary subtraction for \(d=1\)
with multiplicity \(\epsilon\).

It remains to eliminate proper divisors \(m=de\) with \(d,e>1\).
The degree-\(e\) coefficient in \(\log B_j\) has Laurent degree at
most \((j-1)\binom{e+1}2\). Induct on \(m\). In the nonorientable
case the smaller-degree contributions are strictly below \(D_m\)
because \(\binom{d+1}2\binom{e+1}2<\binom{de+1}2\).
In the orientable case the corresponding difference is
\[
 D_m-D_d\binom{e+1}2=
 \frac{de(e-1)}4\big((b_0+s)(d+1)-2b_0\big)>0.
\]
For \(m>1\), the degree-one balanced logarithm in the orientable
subtraction also falls below \(D_m\), since \(s>0\).
What remains in degree \(m\) is \(\sum_jn_{m,j}X_1^{j-1}\).
Since the top Laurent coefficient of \(X_1^D\) is \((-1)^D\),
the two leading coefficients give the asserted signs.
\end{proof}

The orientable formula also recovers \(s\) at degree two. The two
diagrams are transposes, so the disconnected coefficient is
\(2\cos(sv)(X_1X_2)^{b_0}\); its first term after the leading one
distinguishes every \(s\ge0\). At zero mass the bound for \(s>0\)
is attained by the cyclic spectral curve \(y^m=\sigma\) in
\(L_+=aM\), taking the labelling with \(\deg M=s\). For \(m\ge2\)
choose a section \(\sigma\) of \(L_+^m\) with simple zeros
(\(\deg L_+^m\ge2g\)); total ramification over these zeros gives
genus \(1+D_m\). For a nonorientable genus-two target with positive
sign, expansion of Corollary~\ref{cor:cubic} in \(v\), in place of
\(X_1\), gives the connected invariants \(28/3,-25/3,589/180,-61/81\)
in source genera four through seven.

\section{The Euler class of a single spin operator}\label{sec:spin-residues}
A single spin line yields a real Euler class before pairing the two normal directions. This class retains local collision data which disappear from the closed partition function. We compute it first on the finite-map component and then include the stable maps with a contracted elliptic tail.
Let $C$ be a smooth closed complex curve of genus $g\ge2$, and let
$\theta$ be a theta characteristic on $C$, a holomorphic line with an
isomorphism $\theta^2\simeq K_C$. If $f:Y\to C$ is a holomorphic
cover, its pullback spin operator on $L=f^*\theta$ is the real
Cauchy--Riemann operator obtained by adding the metric conjugate of
$df:L\to K_YL^{-1}$ to $\bar\partial_L$. Lee and Parker
\cite{lpob} constructed its real obstruction bundle on the space of
stable maps. Its restriction to unramified covers has rank zero, and
the orientation of a regular cover is its spin parity. Ramification
produces a bundle of positive rank. We determine its Euler class on
the first Hurwitz spaces on which this occurs.

Write $H_m=1+m(g-1)$. In degree $m\ge2$ and source genus $H_m+1$, the
connected unmarked finite maps form a compact smooth stack surface
$\mathcal F_m$. Its general point is a cover with two simple
ramification points. For a connected boundary divisor $B$ of nodal
covers put $s_\theta(B)=(-1)^{h^0(\widetilde f^*\theta)}$, where
$\widetilde f$ is the normalization of a member. For a divisor of
covers with a single triple ramification point $p$, put
$s_\theta(B)=(-1)^{h^0(f^*\theta(p))}$. Both are parities of theta
characteristics. Our divisor formula is
\begin{equation}\label{p8-eq:intro-divisor}
       e(\operatorname{coker}D_\theta)
          =-\frac14\sum_{B} s_\theta(B)[B]
             \quad\text{in }H^2(\mathcal F_m;\QQ).
\end{equation}
Here $B$ runs over the reduced divisors of the two indicated types.
At two simple ramification points over the same value the obstruction
has rational Euler class zero. The equality holds before integrating
against another class.

\subsection{The spin obstruction and the residue pairing}\label{p8-sec:residue}

Let \(f:Y\to C\) be a finite map of smooth compact curves and fix an
isomorphism \(\theta^2\simeq K_C\). Put \(L=f^*\theta\) and
\(P=K_YL^{-1}\). The differential of \(f\) gives a holomorphic map
\(\sigma:L\to P\) whose divisor is the ramification divisor \(R\).
Write \(\deg R=2r\). A Hermitian metric on \(L\) gives a conjugate-linear
identification \(c:P\to\Lambda^{0,1}T^*Y\otimes L\). The real spin
operator and the complex space of ramification residues are
\begin{equation}\label{p8-eq:residue-operator}
 D_\theta=\bar\partial_L+c\sigma,\qquad
 \mathcal E_f=H^0(Y,P/\sigma L).
\end{equation}
The identity
\(\|D_\theta u\|^2=\|\bar\partial u\|^2+\|c\sigma u\|^2\)
follows from Stokes' theorem applied to \(\sigma u^2\). In particular
\(D_\theta\) is injective and its cokernel \(V_\theta\) has real rank
\(2r\). We orient it by continuing the complex determinant orientation
along \(\bar\partial_L+t c\sigma\), \(0\le t\le1\), with the dual
determinant convention for a cokernel \cite[\S12]{lpob}.

\begin{theorem}\label{p8-thm:residue}
Over a family of finite maps with smooth domains and fixed total
ramification \(2r\), the complexification of \(V_\theta\) identifies
with the complex orthogonal bundle \(\mathcal E\) whose form is
\begin{equation}\label{p8-eq:ramification-residues}
 b(v,v')=\sum_{p\in R_{\rm red}}\operatorname{Res}_p
                         \frac{vv'}{\sigma}.
\end{equation}
More precisely, up to the universal normalization of the Serre pairing,
\(V_\theta\) identifies with a positive real form of \(\mathcal E\).
At each map the image \(S=H^0(P)/\sigma H^0(L)\) is a maximal
isotropic subspace. Projection from \(V_\theta\) to
\(\mathcal E/S=S^*\) is a real isomorphism of orientation sign
\((-1)^{h^0(L)}\).
\end{theorem}

\begin{proof}
Resolve \([L\xrightarrow{\sigma}P]\) by its Dolbeault complex. In
degrees zero, one and two the terms are \(\Omega^0(L)\),
\(\Omega^{0,1}(L)\oplus\Omega^0(P)\) and \(\Omega^{0,1}(P)\).
The differentials, written in local coefficients, are
\[
 d_0u=(\bar\partial u,\sigma u),\qquad
 d_1(\beta,v)=\bar\partial_Pv-\sigma\beta .
\]
The sheaf map is injective and its cokernel has finite support.
Hence the middle cohomology is \(\mathcal E_f\) and all other
cohomology vanishes. On the middle term take the symmetric bilinear
pairing
\[
     B\big((\beta,v),(\beta',v')\big)
                      =\int_Y(\beta v'+\beta'v),
\]
with the one-form factors in a fixed order. Integration by parts
identifies the \(B\)-annihilator of \(\operatorname{im}d_0\) with
\(\ker d_1\). On cohomology \(B\) is \eqref{p8-eq:ramification-residues}
up to a constant. Indeed, lift local representatives of \(v\) at the
zeros of \(\sigma\) by a cutoff to a smooth section \(\widetilde v\)
of \(P\). A closed middle cochain representing \(v\) is then
\((\bar\partial\widetilde v/\sigma,\widetilde v)\). Pairing two such
lifts and applying Stokes off disks about the zeros gives the sum of
the boundary integrals of \(vv'/\sigma\). At a zero of order \(d\),
the local form on \(\CC[z]/(z^d)\) takes the coefficient of \(z^{d-1}\)
in the product, after multiplication by a unit. It is nondegenerate,
and distinct support points are orthogonal.

The conjugate-linear involution
\[
                J(\beta,v)=(c(v),c^{-1}(\beta))
\]
makes \(B(\,\cdot\,,J\,\cdot\,)\) a positive Hermitian \(L^2\)
metric after a constant normalization. In a holomorphic frame of
squared norm \(h\), \(J\) sends the local coefficients \((\beta,v)\)
to \((h^{-1}\overline v,h\overline\beta)\). The harmonic middle
space is the Hermitian orthogonal complement of
\(\operatorname{im}d_0+J\operatorname{im}d_0\). In particular it is
\(J\)-invariant. Projection to the first factor identifies the
\(J\)-fixed part of the middle term with \(\Omega^{0,1}(L)\), and
sends \((1+J)d_0u/2\) to \(D_\theta u/2\). A fixed vector real
orthogonal to all these projections is Hermitian orthogonal to
\(\operatorname{im}d_0\), by testing both \(u\) and \(iu\). It is
then also orthogonal to \(J\operatorname{im}d_0\). Thus the real
cokernel identifies with the fixed harmonic space. Elliptic
stabilization gives the identification in families since the harmonic
dimension is constant. Equivalently the orthogonal structure group of
\(\mathcal E\) reduces from \(\mathrm O(2r,\CC)\) to
\(\mathrm O(2r,\RR)\).

The residue theorem makes \(S\) isotropic, and Riemann--Roch gives
\(\dim_\CC S=r\). The long exact sequence identifies
\(\mathcal E_f/S\) with the annihilator of \(\sigma H^0(L)\) in
\(H^1(L)=H^0(P)^*\). On a harmonic class it is the map
\((\beta,v)\mapsto[\beta]\). The image of the zeroth-order term \(c\sigma\) on the Dolbeault kernel is the Hermitian dual of \(\sigma H^0(L)\) in \(H^1(L)\).
Consequently the map from the fixed harmonic space to \(S^*\)
is the cokernel projection of Lee and Parker
\cite[Lemma~12.1 and Proposition~12.2]{lpob}. Its sign is
\((-1)^{h^0(L)}\), from conjugating the complex kernel before taking
its determinant. This last comparison is at a single map and uses
no constant-rank assumption on the Dolbeault cohomology of nearby maps.
\end{proof}

Write \(\kappa=c_1(K_C)\) and use rational cohomology. On a stratum
of fixed ramification profile we may pass to a finite cover ordering
the ramification points. Denote by \(\operatorname{ev}_p\) evaluation
at a marked ramification point \(p\). The following formula concerns
the restriction to this stratum.

\begin{theorem}\label{p8-thm:profiles}
If a ramification index is even, \(e(V_\theta)=0\). If all indices
are odd, write \(R=2D\), \(A=L(D)\), and put
\(I=H^0(Y,L(D)/L)\). Then
\begin{equation}\label{p8-eq:odd-profile-euler}
                 e(V_\theta)=(-1)^{h^0(A)}c_r(I).
\end{equation}
For a fixed target an index \(2s+1\ge5\) also gives zero. If the
ramification consists of \(r\) points of index three, the restriction
to the ordered stratum is
\begin{equation}\label{p8-eq:triple-profile-euler}
             e(V_\theta)=(-1)^{h^0(A)}6^{-r}
                           \prod_{p\in R_{\rm red}}\operatorname{ev}_p^*\kappa .
\end{equation}
\end{theorem}

\begin{proof}
At an even index the corresponding orthogonal residue summand has odd
complex rank. A sum of complex orthogonal bundles admits a positive
real form obtained by taking a positive real form in each summand.
Any two such reductions are isomorphic, since the space of reductions
in each fiber is contractible. After orienting the summands on a
finite cover, the Euler class thus has the Euler class of an
odd-rank oriented real bundle as a factor. The latter vanishes over
\(\QQ\). Transfer gives the assertion on the original stratum.

For odd indices \(A^2=K_Y\) by Riemann--Hurwitz. Multiplication by
\(\sigma\) identifies \(P\) with \(L(2D)\). In this description \(I\)
is a maximal isotropic subbundle of \(\mathcal E\), since the local
meromorphic representatives in \(I\) have product with no residue
in \eqref{p8-eq:ramification-residues}. The two maximal isotropics have
\[
                     k=\dim(S\cap I)=h^0(A)-h^0(L).
\]
This is the residue construction for a theta characteristic
\cite[\S1]{mumford}. The line \(A\) is the theta characteristic of the
oddly ramified cover used in spin Hurwitz theory \cite[\S1]{lp}.
For two maximal isotropics meeting in dimension \(k\), the real
projections from a positive real form to their duals induce
orientations which differ by \((-1)^{r-k}\). To see this, move the
pair through maximal isotropics with intersection of dimension \(k\)
to a pair in hyperbolic coordinate form. In the \(r-k\) exchanged
coordinates the projections differ by complex conjugation. Applying
Theorem~\ref{p8-thm:residue} and \(c_r(I^*)=(-1)^r c_r(I)\) gives the
sign \((-1)^{h^0(L)+r-k+r}=(-1)^{h^0(A)}\). The bundle \(I\) exists
over the entire profile stratum. The parity of the theta
characteristic \(A\) is locally constant, so the pointwise orientation
comparison proves \eqref{p8-eq:odd-profile-euler} across jumps of \(h^0(L)\).

At an index \(2s+1\) point with image \(x\), the filtration by pole
order on \(I\) has graded lines
\(\theta_x\otimes(T_pY)^{\otimes j}\), \(1\le j\le s\). The
leading nonzero jet of \(f\) identifies \((T_pY)^{\otimes(2s+1)}\)
with \(T_xC\). The first Chern classes of the graded lines are
therefore
\[
              \left(\frac12-\frac{j}{2s+1}\right)
                         \operatorname{ev}_p^*\kappa .
\]
When \(s\ge2\), two such factors already have product zero since
\(\kappa^2=0\) on \(C\). For \(s=1\) the factor is
\(\operatorname{ev}_p^*\kappa/6\), giving
\eqref{p8-eq:triple-profile-euler}.
\end{proof}

On the stratum with \(r\ge2\) distinct triple critical values,
the product in \eqref{p8-eq:triple-profile-euler} itself vanishes in
ordinary cohomology. It pulls back from the top cohomology of the
noncompact ordered configuration space of \(r\) points on \(C\).
A prescribed coincidence of two values gives zero by \(\kappa^2=0\).
The class of a family meeting the boundary can nevertheless have a
contribution supported there.

\begin{corollary}\label{p8-cor:nodal-spin-matching}
On a stratum of finite nodal maps with fixed dual graph, label the
components, the nodes, and the two preimages of each node on a finite
cover. Let \(V_{\theta,i}\) be the oriented cokernel on the \(i\)-th
normalized component and let \(\operatorname{ev}_j\) send a map to
the image \(x_j\) of its \(j\)-th node.
Then the restriction of the nodal Euler class is
\[
             \prod_i e(V_{\theta,i})\,
             \prod_j\operatorname{ev}_j^*c_1(\theta).
\]
An unramified component contributes its index-zero determinant sign.
\end{corollary}

\begin{proof}
In the \(W^{1,p}\) domain with \(p>2\), the nodal domain is the
kernel of the surjection taking the differences of values over the
nodes, with target \(\bigoplus_j\theta_{x_j}\). The spin operators
on the normalized components are injective. The exact sequence of
kernels and cokernels therefore gives
\[
 0\longrightarrow\bigoplus_j\theta_{x_j}\longrightarrow
 V_{\theta,\mathrm{nod}}\longrightarrow\bigoplus_iV_{\theta,i}
 \longrightarrow0.
\]
The determinant isomorphism gives the matching lines their complex
orientations, by continuation from the complex-linear normalization
sequence \cite{zinger}. All ranks in the sequence are even, so
reordering the factors preserves the product orientation. Interchange
of the two branches at a node tensors its matching line with a flat
sign line and leaves its rational first Chern class unchanged.
\end{proof}

\section{Collision divisors and elliptic tails}\label{p8-sec:hurwitz}

Fix a smooth compact curve \(C\) of genus \(g\ge 2\), put
\(H_m=1+m(g-1)\), and let \(m\ge 2\). Inside the stack of connected
stable maps \(\overline{\mathcal M}_{H_m+1}(C,m)\), let
\(\mathcal F_m=\mathcal F_m(C)\) be the locus of maps finite over \(C\).
These maps are flat, since the source is a nodal curve and its
structure sheaf is torsion-free over the smooth target. The locus is
open by properness of the universal curve and openness of
quasi-finiteness. Its complement is a union of open and closed
elliptic-tail loci, as proved in Section~\ref{p8-sec:tails}. Thus
\(\mathcal F_m\) is proper. Bryan and Pandharipande introduced this
surface and its branch morphism in their calculation of local
invariants of curves \cite[\S\S1.3.2, 4]{bp01}.

For a finite map, Riemann--Hurwitz on the normalization and the genus
formula for a nodal curve give
\[
           2=\deg R+2\,\#\{\text{nodes}\},
\]
where \(R\) is the ramification divisor on the normalization. Thus a
nodal map has one node and its normalization is \'etale. At a smooth
ramification point of index \(b\) the local presentation \(w=z^b\)
identifies the relative cotangent complex with the module
\(\mathcal O/(z^{b-1})\) in degree zero. At the node there are
coordinates \(uv=0\), \(w=u+v\), since both branches are unramified.
Writing its local algebra over the target as
\(\mathcal O_C[u]/(u^2-wu)\) gives the relative cotangent complex
\(\mathcal O/(2u-w)\) in degree zero. Here \(2u-w=u-v\) is a
nonzerodivisor and the quotient has length two. In each case the
module has projective dimension one. Its sheaf \(\mathcal{H}\!om\)
to \(\mathcal O\) vanishes, and its sheaf \(\mathcal{E}\!xt^1\) has the
same finite length as the module. The local-to-global sequence for
deformations over \(C\) consequently gives tangent dimension two and
zero obstruction. This proves smoothness of the stack and identifies
its virtual class with its fundamental class.

The branch morphism
\(\beta:\mathcal F_m\to C^{(2)}\) sends a smooth cover to its branch
divisor and a nodal one to twice the image of the node. Let
\(\Delta=\{2x:x\in C\}\). The connected components of the reduced
preimage of \(\Delta\) are of three kinds.
\begin{itemize}
\item \(N_i\) consists of nodal covers. The two points above the node
  belong to a possibly disconnected \'etale normalization.
\item \(T_j\) consists of smooth covers having a single ramification
  point of index three.
\item \(U_k\) consists of smooth covers having two distinct simple
  ramification points over the same point of \(C\).
\end{itemize}
For any such component \(B\), write \(a_B:B\to C\) for its branch
value and \(d_B=\int_B a_B^*[{\rm pt}]\). This is a stack degree.
For a fixed branch value it sums \(1/|\operatorname{Aut}_C(f)|\)
over the corresponding isomorphism classes. Indeed, the monodromy
data over the complement of the branch value are discrete and finite
in fixed degree, and move uniquely up to isomorphism in a disk of
branch values. In the nodal case the data consist of an \'etale cover
and an unordered pair of distinct sheets to be identified, with
connected quotient. Automorphisms interchanging the branches are
included.

In a versal deformation of a nodal map the local model is
\(uv=t,\ w=u+v+v_0\). Its two critical values have squared difference
\(16t\). At a point of \(T_j\) the model is \(w=z^3+uz+v\), for
which the squared difference is a nonzero constant times \(u^3\).
At a point of \(U_k\) the two critical values \(v_1,v_2\) are
independent parameters on an ordering cover; their squared
difference is \((v_1-v_2)^2\). The local deformation calculation
above gives all the parameters in these models. In particular the
reduced divisors are smooth and pairwise disjoint, and
\begin{equation}\label{p8-eq:hurwitz-pullback}
        \beta^*\Delta=\sum_i N_i+3\sum_j T_j+2\sum_k U_k.
\end{equation}
The multiplicities are orders on stack charts. Since the pullback of
\(\Delta\) to \(C\times C\) is twice the diagonal, \(\Delta^2=4-4g\).
Restricting \eqref{p8-eq:hurwitz-pullback} to a component of multiplicity
\(\ell_B\) gives
\begin{equation}\label{p8-eq:hurwitz-intersection}
                         B^2=\frac{4-4g}{\ell_B}\,d_B .
\end{equation}
These are the branch intersections used in \cite[\S4.3]{bp01}.

For a theta characteristic \(\theta\) define the following signs on the
reduced divisors. On \(N_i\), let \(\widetilde f\) be the normalization
of the cover. On \(T_j\), write \(p\) for the ramification point. Put
\begin{equation}\label{p8-eq:hurwitz-signs}
 s_\theta(N_i)=(-1)^{h^0(\widetilde f^*\theta)},\qquad
 s_\theta(T_j)=(-1)^{h^0(f^*\theta(p))}.
\end{equation}
In each expression the line is a theta characteristic on a smooth
curve, and its parity is locally constant. On a disconnected
normalization \(h^0\) is the sum over its components.

\begin{theorem}\label{p8-thm:boundary}
In \(H^2(\mathcal F_m;\QQ)\) the Euler class of the oriented spin
obstruction plane is
\begin{equation}\label{p8-eq:hurwitz-euler}
 e(V_\theta)=-\frac14
    \left(\sum_i s_\theta(N_i)[N_i]+\sum_j s_\theta(T_j)[T_j]\right),
\end{equation}
where brackets denote the classes of the reduced divisors.
For \(D=D_\theta\oplus D_\eta\), the finite-map contribution is
\begin{equation}\label{p8-eq:hurwitz-finite}
 F_m(D)=-\frac{g-1}{4}\left(
    \sum_i d_{N_i}s_\theta(N_i)s_\eta(N_i)
    +\frac13\sum_j d_{T_j}s_\theta(T_j)s_\eta(T_j)\right).
\end{equation}
\end{theorem}

\begin{proof}
The cokernels over finite nodal maps extend the oriented plane on the
smooth locus by the normal obstruction bundle construction
\cite[Theorem~8.1 and Lemma~11.1]{lpob}. Off the \(N_i\) and \(T_j\) there are two
distinct simple ramification points. By Theorem~\ref{p8-thm:residue}
their residue spaces are orthogonal lines. Ordering the points and
choosing a unit vector for each complex quadratic form gives a finite
cover trivializing the orthogonal bundle on this entire open set,
including the \(U_k\). Uniqueness up to homotopy of its positive
real form and transfer imply that the rational Euler class restricts
to zero there. By the Thom isomorphism for the disjoint smooth
divisors, it is therefore a linear combination of the classes
\([N_i]\) and \([T_j]\).

Set \(\kappa=c_1(K_C)\). On \(N_i\) the normalization sequence of
Corollary~\ref{p8-cor:nodal-spin-matching} gives
\(e(V_\theta)|_{N_i}=s_\theta(N_i)a_{N_i}^*\kappa/2\).
On \(T_j\) Theorem~\ref{p8-thm:profiles} gives
\(e(V_\theta)|_{T_j}=s_\theta(T_j)a_{T_j}^*\kappa/6\).
Integrating over a component \(B\) and dividing by its nonzero
self-intersection \eqref{p8-eq:hurwitz-intersection} determines its
coefficient. It is \(-s_\theta(B)/4\) in both cases. Finally the
oriented cokernel for \(D\) is the direct sum of the two spin
obstruction planes. Integrating the product of their Euler classes
and applying \eqref{p8-eq:hurwitz-intersection} gives
\eqref{p8-eq:hurwitz-finite}.
\end{proof}
\subsection{Contracted elliptic tails}\label{p8-sec:tails}

Let \(N\to C\) be a complex vector bundle of rank two and degree
\(2g-2\), and \(D\) a real Cauchy--Riemann operator on \(N\). Assume
that the pullback operator on every smooth connected cover of \(C\)
is injective. On an \'etale cover the index is zero, so the operator
is invertible. Orient all normal determinants by continuation from
complex-linear operators, and denote the determinant sign of an
\'etale cover \(f\) by \(\sigma_D(f)\). The signed stack degree of
the connected \'etale covers of degree \(m\) is
\[
 c_m(D)=\sum_{[f:Y\to C]}
                  \frac{\sigma_D(f)}{|\operatorname{Aut}_C(f)|}.
\]
Write \(L_{m,H_m+1}(D)\) for the absolute normal Euler number on
\(\overline{\mathcal M}_{H_m+1}(C,m)\) and \(F_m(D)\) for its
integral over the finite-map locus. Put \(F_1(D)=0\), since a
finite degree-one cover is an isomorphism. When \(D\) is the normal
operator of a super-rigid embedded curve, the fixed-target
comparison identifies this absolute Euler pairing with the local
Gromov--Witten degree \cite[Theorem~1.2(2)--(3)]{zinger}.
For $m\ge2$ the relatively oriented Euler number of
\cite[Definition~6.2 and Remark~6.4]{bai} differs from this absolute
number by the determinant sign of the operator on $C$.


\begin{proposition}\label{p8-prop:tails}
The complement of the finite-map locus is a disjoint union of
open and closed components indexed by the connected degree-\(m\)
\'etale covers \(f:Y\to C\). Each component has its ordinary
fundamental class as virtual class. Their contributions give
\begin{equation}\label{p8-eq:tails-identity}
         L_{m,H_m+1}(D)=F_m(D)-\frac{m(g-1)}{12}\,c_m(D).
\end{equation}
\end{proposition}

\begin{proof}
Consider a stable map, and let \(s\) be the number of nonconstant
components on the normalization of its domain. Write \(R\) for the
sum of their ramification divisors, \(g_0\) for the sum of the
genera of the normalized contracted components, and \(\Gamma\) for
the connected dual graph. Riemann--Hurwitz and the arithmetic genus
give
\[
             2=s+\frac{\deg R}{2}+g_0+b_1(\Gamma).
\]
If \(s=2\), both positive-degree components are \'etale and the
graph is a tree. A contracted rational vertex must have valence at
least three, which is impossible in a tree with at most two leaves.
If \(s=1\) and \(\deg R=2\), the same argument excludes a contracted
subtree. If \(s=1\) and \(\deg R=0\), the nonconstant component has
smooth normalization \(Y\) of genus \(H_m\), and \(g_0+b_1(\Gamma)=1\).
When \(g_0=1\) the tree consists of \(Y\) joined directly to a
smooth genus-one component. When \(b_1(\Gamma)=1\), every contracted
vertex has valence at least three. The identity
\(\sum_v(\operatorname{val}(v)-2)=0\) leaves just a self-node on
\(Y\), or one contracted rational vertex with a self-loop joined
to \(Y\) by one edge. Thus every map with a contracted component
consists of a smooth connected \'etale cover with a one-pointed
stable curve of genus one attached.

For fixed \(f\), gluing at \(x\in Y\) defines a proper morphism from
\[
 \mathcal G_f=
  \big[(Y\times\overline{\mathcal M}_{1,1})/\operatorname{Aut}_C(f)\big]
      \longrightarrow\overline{\mathcal M}_{H_m+1}(C,m).
\]
The deck group acts on \(Y\), and the pointed tail retains its stack
stabilizers. To compute the deformation theory at
\(Y\cup_{x=y}E\), first preserve the attaching node. The cover is
rigid because \(df:T_Y\to f^*T_C\) is an isomorphism. Its marked
point moves freely. The tail deforms as a pointed stable curve with
constant image in \(C\). Varying this image makes the matching
equation at the node transverse. The remaining obstruction is
\(T_{f(x)}C\otimes H^1(E,\mathcal O_E)\). The boundary map from
the smoothing line is
\begin{equation}\label{p8-eq:tails-smoothing}
 T_xY\otimes T_yE\ \xrightarrow{\,df_x\otimes\delta_y\,}\
 T_{f(x)}C\otimes H^1(E,\mathcal O_E),
\end{equation}
up to the choice of sign for the connecting homomorphism. Indeed,
in the plumbing \(uv=t\), the linear term on the cover gives the
principal part \(df_x(t/v)\) on the tail. Thus \(\delta_y\) is the
connecting homomorphism for
\(0\to\mathcal O_E\to\mathcal O_E(y)\to T_yE\to0\). The dualizing
sheaf is trivial for both a smooth elliptic curve and an irreducible
one-nodal genus-one curve. Serre duality and Riemann--Roch give
\(H^1(\mathcal O_E(y))=0\) and \(H^0(\mathcal O_E(y))=\CC\).
Hence \eqref{p8-eq:tails-smoothing} is an isomorphism on all of
\(\mathcal G_f\).

The full stable-map theory is therefore unobstructed of complex
dimension two along \(\mathcal G_f\), and its tangent space consists
of the attachment and pointed-tail deformations. These are precisely
the tangents to the source of the gluing morphism. Gluing identifies
the stabilizers and is injective on geometric points, since the
positive-degree component and the attaching node are intrinsic. It
is thus a proper representable \'etale open immersion, and its image
is closed as well. This proves the asserted component description
and the fundamental-class claim. It also gives the properness used in
Section~\ref{p8-sec:hurwitz}. An alternative description by lifting to
the maximal \'etale cover appears in \cite[\S2.2]{bp01}.

It remains to orient the normal obstruction. On \(Y\) the operator
\(f^*D\) is invertible. On the constant tail the normal operator
is \(\bar\partial\) with coefficients in \(N_{f(x)}\). Its kernel
consists of the constants, which evaluate isomorphically onto the
matching fiber. The normalization exact sequence at the attaching
node consequently identifies the full normal cokernel with
\[
                    H^1(E,\mathcal O_E)\otimes N_{f(x)}.
\]
This uses the arithmetic \(H^1\) also for a nodal tail. In families
the first factor is the dual of the Hodge line \(\mathbb E\). The
normalization isomorphism of determinants gives the constants and the
matching fiber their complex orientations; their determinants cancel.
Continuing from the complex-linear normalization sequence leaves
the factor \(\sigma_D(f)\) from the invertible cover operator
\cite[\S2]{zinger}, \cite[Appendix~B]{bai}. Thus the full normal
orientation is the complex orientation of
\(\mathbb E^\vee\otimes f^*N\) multiplied by \(\sigma_D(f)\).

Write \(\lambda=c_1(\mathbb E)\). On the product of the two curves
\[
        c_2(\mathbb E^\vee\otimes f^*N)=-\lambda\,c_1(f^*N).
\]
The genus-one discriminant has a simple zero at the nodal boundary
and weight twelve, so \(12\lambda\) is the boundary class on
\(\overline{\mathcal M}_{1,1}\). The boundary point has stabilizer
of order two, whence \(\int\lambda=1/24\). Finally
\(\int_Y c_1(f^*N)=m(2g-2)\). Dividing by
\(|\operatorname{Aut}_C(f)|\) and summing proves
\eqref{p8-eq:tails-identity}.
\end{proof}

We record the implication for the exponents of the elementary
factors. Let \(n_{d,a}\) be the coefficients in the expansion of
the connected local series in the connected elementary series
\(G_a(q^d,t)\), with \(a\ge H_d\). Here $G_a$ is the
connected Gromov--Witten series of the elementary genus-$a$ curve,
the logarithm of its normalized partition function. A connected
degree-$j$, genus-$h$ contribution has monomial $q^jt^{2h-2}$.
The same counts are attached to the balanced universal factors
in Section~\ref{p8-sec:exponents}. Write \(C_j(a)\) for the
unsigned connected \'etale degree-\(j\) count of a genus-\(a\)
curve, and \(F_j^{\rm elem}(a)\) for the finite-map Euler number of
its doubled spin operator at first excess. All \'etale signs of this
operator are positive. The elementary series begins in degree \(j\)
at genus \(1+j(a-1)\), and its degree-one leading coefficient is one
\cite[Propositions~3.2--3.3]{ip}, \cite[Proposition~3.46]{diw}. At genus
\(H_m\), comparison therefore gives
\[
               c_m(D)=\sum_{d\mid m} n_{d,H_d} C_{m/d}(H_d).
\]
At genus \(H_m+1\), a proper divisor \(d\mid m\) can still contribute
only the factor of genus \(H_d\), since increasing that genus by one
raises the least source genus by \(m/d\). For the equal-degree
factor of genus \(H_m\) the finite term is \(F_1^{\rm elem}(H_m)=0\).
In each divisor term the tail multiplier in
\eqref{p8-eq:tails-identity} is the same, since
\((m/d)(H_d-1)=m(g-1)\). The least-genus identity cancels their sum
against the tail contribution on the left and yields
\begin{equation}\label{p8-eq:tails-exponent}
          n_{m,H_m+1}=F_m(D)-
          \sum_{\substack{d\mid m\\d<m}}
                   n_{d,H_d}\,F_{m/d}^{\rm elem}(H_d).
\end{equation}

\section{Signed Hurwitz numbers and their elementary multiplicities}\label{p8-sec:characters}

Let $a,b$ be the theta characteristics of the two normal lines,
$w=ba^{-1}\in H^1(C;\ZZ/2)$ their difference, and
$\epsilon=(-1)^{h^0(a)+h^0(b)}$. Put $H_m=1+m(g-1)$ and
$c_m=L_{m,H_m}$. For a cell $u$ of a partition $\lambda$, let
$h(u)$ denote its hook length and $c(u)$ its column minus its row.
Write $H_\lambda=\prod_{u\in\lambda}h(u)$. Let $\mathcal P_w$
consist of all partitions if $w=0$ and all self-conjugate partitions
if $w\ne0$, including the empty partition in either case. Set
\[
 Z_w(y)=\sum_{\lambda\in\mathcal P_w}
               H_\lambda^{\,2g-2}y^{|\lambda|}.
\]

\begin{theorem}\label{p8-thm:hooks}
The connected local counts at the first two possible source genera
satisfy $c_m=\epsilon^m[y^m]\log Z_w(y)$ and
\begin{equation}
 L_{m,H_m+1}
 =-\frac{\epsilon^m(g-1)}{12}[y^m]\,
  \frac{\displaystyle\sum_{\lambda\in\mathcal P_w}
         H_\lambda^{\,2g-2}\!
         \left(\sum_{u\in\lambda}h(u)^2\right)y^{|\lambda|}}
       {Z_w(y)}.
 \label{p8-eq:character-first-excess}
\end{equation}
\end{theorem}

For equal spin structures the right side is the first correction
in the balanced hook formula. The restriction to self-conjugate
partitions comes from a sign on one handle of the target.

\subsubsection{The spin sign of a collision}
At a nodal cover let $f:Y\to C$ be its \'etale normalization,
including all components. At a smooth cover with one triple
ramification point let $f:Y\to C$ be the cover itself. Write
$R$ for the ramification divisor of $f$, so $R=0$ in the first
case and $R=2r$ in the second. The two signs in the boundary
intersection are the parities of
\[
       f^*a\otimes\mathcal O_Y(R/2),\qquad
       f^*b\otimes\mathcal O_Y(R/2).
\]
For an odd branch these are the spin structures of spin Hurwitz
theory \cite{eop}. Their product has a simpler expression.

\begin{lemma}\label{p8-lem:odd-branch-sign}
Let $\alpha\in H^1(C;\ZZ/2)$ be the sign character of the sheet
permutations of $f$. Then the product of the two spin signs equals
\begin{equation}
       \epsilon^m(-1)^{\langle w\smile\alpha,[C]\rangle}.
 \label{p8-eq:odd-branch-sign}
\end{equation}
\end{lemma}
\begin{proof}
The bundle $E=f_*\mathcal O_Y(R/2)$ has rank $m$ and carries the
orthogonal form $\operatorname{Tr}(uv)$. Around $x=z^3$ its ramified
summand has basis $z^{-1},1,z$ over $\CC\{x\}$. The matrix of the
trace form in this basis has antidiagonal entries $3$ and all
other entries zero. The remaining sheets have the ordinary
permutation form. Thus the pairing extends nondegenerately over the
branch value. On its complement the determinant character of $E$
is the permutation sign. A 3-cycle is even, so this character
extends across the missing point and $w_1(E)=\alpha$.

For an orthogonal bundle of rank $m$ over a spin curve the
mod-two index formula is
\[
 h^0(a\otimes E)\equiv
 m h^0(a)+q_a(w_1(E))+\langle w_2(E),[C]\rangle\pmod2,
\]
where $q_a(v)=h^0(a\otimes v)-h^0(a)\pmod2$
\cite{ati,baraglia}. By the projection formula the left side is
the parity on $Y$. Adding this identity for $a$ and $b$ cancels
$w_2(E)$. Since $q_b(v)-q_a(v)=\langle w\smile v,[C]\rangle$,
the result follows. The same calculation applies to a disconnected
source by direct sum.
\end{proof}

\subsubsection{One signed handle and one marked collision}
We remove the common factor $\epsilon^m$ in degree $m$ until the
end of the calculation. For $w\ne0$ choose a simple nonseparating
circle representing its Poincar\'e dual as the first generator in a
surface presentation, disjoint from the branch value. The remaining
sign in~\eqref{p8-eq:odd-branch-sign} is $\operatorname{sgn}(A_1)$
for its sheet permutation $A_1$. Take $\chi=\operatorname{sgn}$ in
this case and $\chi=1$ for $w=0$. For $G=S_m$ put
\[
 K_\chi=\sum_{A,B\in G}\chi(A)\,ABA^{-1}B^{-1}\in\CC[G].
\]
On an irreducible representation $V$ of dimension $d$ and character
$\psi$, this central element has scalar
\[
 \frac{|G|}{d^2}\sum_{A\in G}\chi(A)|\psi(A)|^2
  =\left(\frac{|G|}{d}\right)^2
       {\bf1}_{\,V\otimes\chi\simeq V}.
\]
Indeed the sum over $B$ of the conjugates of $A^{-1}$ acts by
$|G|\overline{\psi(A)}/d$, and the last equality is character
orthogonality. The other $g-1$ handles contribute $K_1$.

The identity coefficient of $K_\chi K_1^{g-1}$ counts the labelled
\'etale covers with this sign. It equals the regular trace divided
by $|G|$. Relabelling sheets divides by $|G|$ once more, giving
the inverse automorphism weights. The summand of $V$ is therefore
$(|G|/d)^{2g-2}$ when $V\otimes\chi\simeq V$. For
$V=V_\lambda$ the hook-length formula gives $|G|/d=H_\lambda$,
and tensoring with sign transposes $\lambda$. The resulting
disconnected series is $Z_w$. Decomposition into connected
components gives the formula for $c_m$, as in \cite[Appendix~F]{WueUnramified}.

Denote by $N_m$ the signed mass over one fixed target point of
\'etale covers with two distinct sheets identified to give a
connected nodal cover. Denote by $T_m$ the signed mass of connected
covers with a single 3-cycle over that point. Both masses still omit
$\epsilon^m$. Before imposing connectivity the first mark has
$\binom m2$ choices on each labelled fiber. After gluing, the marked
component is unique, and all its companions are unmarked \'etale
covers. Removing their generating function gives
\begin{equation}
    \sum_{m\ge1}N_m y^m=\frac{y^2 Z_w''(y)}{2Z_w(y)}.
 \label{p8-eq:character-nodal}
\end{equation}
In particular this includes a node joining two distinct components
of the normalization. For the triple mark the handle product is
multiplied by the class sum of all 3-cycles. The Jucys--Murphy
elements $J_j=\sum_{i<j}(ij)$ satisfy
\[
     \sum_j J_j^2=\binom m2\,1+\sum_{\tau\text{ a 3-cycle}}\tau .
\]
Indeed the nondiagonal terms of $J_j^2$ list the 3-cycles with
largest letter $j$, each once. In the Young basis the $J_j$ act
by the contents of the cells containing $j$ \cite{jucys}. Removing
the unmarked components as above now gives
\begin{equation}
 \sum_{m\ge1}T_m y^m
  =\frac{1}{Z_w(y)}\sum_{\lambda\in\mathcal P_w}
      H_\lambda^{\,2g-2}
      \left(\sum_{u\in\lambda}c(u)^2-\binom{|\lambda|}{2}\right)
      y^{|\lambda|}.
 \label{p8-eq:character-triple}
\end{equation}

\begin{proof}[Proof of Theorem~\ref{p8-thm:hooks}]
The two collision intersections give the finite contribution
$-\epsilon^m(g-1)(3N_m+T_m)/12$. The contracted elliptic tails add
$-m(g-1)c_m/12$. Since $y(\log Z_w)'=y Z_w'/Z_w$, the three
terms over the common denominator $Z_w$ have total insertion on
a partition of size $m$ equal to
\[
 3\binom m2+\sum_{u\in\lambda}c(u)^2-\binom m2+m
 =m^2+\sum_{u\in\lambda}c(u)^2
 =\sum_{u\in\lambda}h(u)^2 .
\]
The last equality is the hook-content square identity \cite{menon}.
Restoring $\epsilon^m$ proves~\eqref{p8-eq:character-first-excess}.
\end{proof}

\subsubsection{Degree three}
Write $s=2g-2$. In degree three the partitions $(3),(2,1),(1^3)$
have hook products $6,3,6$, respectively. Only $(2,1)$ is
self-conjugate. Equations~\eqref{p8-eq:character-nodal}
and~\eqref{p8-eq:character-triple}, including the lower coefficients in
the division by $Z_w$, give
\[
\begin{array}{c|cc}
 &N_3&T_3\\ \hline
 w=0   &6\cdot6^s+3\cdot3^s-2\cdot2^s &(4^g-1)3^s\\
 w\ne0 &3^{2g-1}&-3^s .
\end{array}
\]
Thus the finite contribution for equal spins is
$-(g-1)(22\cdot6^s+8\cdot3^s-6\cdot2^s)/12$, and for distinct
spins it is $-2\epsilon(g-1)3^s/3$. At this prime degree the
proper divisor one contributes $\epsilon$ times the finite degree
of the balanced genus-$g$ core. The degree-one finite contribution
at positive excess is zero. Cancellation of the elliptic tails
therefore gives, for $a\ne b$,
\begin{equation}
        n_{3,H_3+1}
       =\frac{\epsilon(g-1)}{6}
          \left(11\cdot6^{\,2g-2}-3\cdot2^{\,2g-2}\right).
 \label{p8-eq:cubic}
\end{equation}
For $g=2$ the ordinary collision masses are $235$ and $135$,
with finite contribution $-70$. For distinct spins they are $27$
and $-9$ before restoring $\epsilon$, with finite contribution
$-6\epsilon$. Their difference gives $n_{3,5}=64\epsilon$.
\subsection{Elementary exponents and unramified quotients}\label{p8-sec:exponents}

Let $F_h(q)$ be the balanced antidiagonal local factor of genus $h$ in
the rational universal algebra of complex one-cycles of virtual dimension
zero \cite{pardon,WueProduct}. We use the convention that $F_h$ has
constant term one. Its image in the algebra of almost-complex cycle
families is the elementary cluster of genus $h$ of Ionel and
Parker \cite{ip,WueProduct}. For the isolated spin germ, the characteristic
insertions at the compact cycles $m[C]$ have a unique balanced lift
\cite[Proposition~3.1]{WueUnramified}
\begin{equation}\label{p8-eq:lifted-product}
  \widetilde P=\prod_{d\ge1,\ h\ge H_d}F_h(q^d)^{n_{d,h}}.
\end{equation}
The exponents are integral and have finite support in $h$ at fixed
$d$. This uses geometric decomposition after compact realization and
the multiplicity-one detection theorem of \cite{WueProduct}. Along the
fixed-target cluster decompositions, connected Gromov--Witten degrees
add on disjoint targets. Thus the same exponents give the expansion
of the connected series in the elementary logarithms. The use of
fixed-target degrees in this expansion relies on that geometric
comparison. In particular no evaluation of arbitrary analytic
parameter families is required here.

Put $n_d=n_{d,H_d}$. In this notation the tail cancellation of
\eqref{p8-eq:tails-exponent} subtracts the equal-spin finite-map integrals
on unramified covers. We next express the subtraction in terms of
the individual maps to $C$.

\subsubsection{A count using the maximal unramified quotient}
Fix a target point $x$. Let $\mathcal B_m(x)$ denote the isomorphism
classes of the finite maps over the nodal and triple divisors with
collision value $x$. Such an object has one distinguished node or one
triple ramification point. Write $k(f)=1$ in the nodal case and
$k(f)=3$ in the triple case, and write $\tau(f)$ for the product of
the two spin parities in~\eqref{p8-eq:intro-divisor}. All automorphisms
in this section fix the target. Put $\Gamma=\pi_1(C)$ and regard
$H(f)=f_*\pi_1(Y)$ as a conjugacy class of subgroups. This is the factorization by the image of the fundamental group used
by Bryan and Pandharipande \cite[\S2.2, Proposition~2.3]{bp01}.
It uses
the connected nodal space $Y$ itself when $f$ has a node.
The connected cover of $C$ corresponding to $H(f)$ is finite:
$f$ lifts to it, and the lift is holomorphic locally on every
component. A nonconstant holomorphic map from a compact component
to a connected Riemann surface forces that surface to be compact.
Composition of degrees shows that $[\Gamma:H(f)]$ divides $m$.
This cover is the maximal unramified quotient of $f$, in the sense
that every factorization through an unramified connected cover
factors through it after making compatible choices of base points.
For finite-index $H\le\Gamma$ write $\sigma(H)$ for the determinant
sign of the spin pair on its unramified cover.

\begin{theorem}\label{p8-thm:transfer}
For every $m\ge1$,
\[
 n_{m,H_m+1}
  =\frac{g-1}{4}\!
     \sum_{f\in\mathcal B_m(x)}
     \frac{\sigma(H(f))-\tau(f)}
          {k(f)\,|\operatorname{Aut}_C(f)|}.
\]
In particular the difference of the exponents for two spin pairs
is obtained by taking the corresponding difference in these
boundary and unramified signs on the same covers.
\end{theorem}
\begin{proof}
By the orbit formula for the determinant signs \cite[Theorem~16.1]{WueUnramified}, there are
integers $b_K$ indexed by conjugacy classes of finite-index subgroups
such that
\[
 \sigma(H)=\sum_{[K]} b_K |(\Gamma/K)^H|,
 \qquad n_d=\sum_{[\Gamma:K]=d} b_K .
\]
In the second sum each conjugacy class is used once. The first formula
is the fixed-point count, or mark, of a unit in the completed Burnside
ring; only subgroups of index dividing $[\Gamma:H]$ contribute.
For a fixed cover $z:Z\to C$ corresponding to $K$ of index $d$,
the number of lifts of $f$ to $Z$ is $|(\Gamma/K)^{H(f)}|$.
The groupoid of such lifts over an object $f$ has mass this number
divided by $|\operatorname{Aut}_C(f)|$, by orbit-stabilizer.
Here $Z$ is a fixed target, so its deck group does not divide this
number. Counting the same lifted maps on $Z$ allows the collision
value to be any of the $d$ preimages of $x$. The mass at each such
value is the same, by continuation along the connected curve $Z$.
An unramified change of target preserves the node or the triple
point and hence the weight $1/k(f)$. All boundary signs in the
equal-spin elementary pair on $Z$ are positive. By
Theorem~\ref{p8-thm:boundary} and $H_d-1=d(g-1)$ we obtain
\[
 F_{m/d}^{\rm elem}(H_d)=
 -\frac{g-1}{4}
  \sum_{f\in\mathcal B_m(x)}
  \frac{|(\Gamma/K)^{H(f)}|}
       {k(f)\,|\operatorname{Aut}_C(f)|}
\]
whenever $d\mid m$. For $d\nmid m$ the fixed sets are empty.
Multiply by $b_K$, sum, and apply~\eqref{p8-eq:tails-exponent} and
the signed boundary formula for $D$.
\end{proof}

\begin{corollary}\label{p8-cor:prime-growth}
Let $a\ne b$ and let $p$ be an odd prime. The signs of $n_{p,H_p}$
and $n_{p,H_p+1}$ are $-\epsilon$ and $\epsilon$, respectively, and
\[
 |n_{p,H_p+1}|\ > \frac{g-1}{4}\binom{p}{2}\,|n_{p,H_p}| .
\]
For each $g$ there is a path of tame structures on a closed
symplectic six-torus for which the isolated positive-sign core at
its two ends has different local elementary exponents at both these
source genera for every odd prime $p$.
\end{corollary}
\begin{proof}
Every collision in prime degree has $H(f)=\Gamma$. Otherwise its
maximal unramified quotient would have degree $p$ and the lift degree
one. A finite flat degree-one cover of a smooth curve is an
isomorphism, inconsistent with either collision. Thus its unramified
sign in Theorem~\ref{p8-thm:transfer} is $\epsilon$. The boundary sign
is $\epsilon^p$ times the permutation sign around a loop Poincar\'e
dual to the difference $ba^{-1}$, by the parity formula in
Section~\ref{p8-sec:characters}. Consequently each negative permutation
sign contributes with sign $\epsilon$. Let $U_-$ be the inverse-deck
mass of connected unramified covers with that negative sign. Each
one supplies $\binom p2$ unordered pairs in the fiber at $x$ whose
identification forms a node of the same character. Their total mass
is $\binom p2 U_-$. In least genus $n_{p,H_p}=-2\epsilon U_-$ by
the unramified formula. Also $U_->0$: choose surface generators
with the chosen loop first, send it to a transposition, send the
generator of another handle to a $p$-cycle, and send all remaining
generators to the identity. The source is connected and the surface
relation holds. These nodes alone give the weak inequality. There is also a nodal cover
with disconnected normalization of degrees $1$ and $p-1$. Choose a
connected cyclic $(p-1)$-cover whose sign on the chosen loop is negative
by sending the loop to a full cycle. Join one of its sheets to the
degree-one sheet at $x$. The resulting map has the same negative
permutation sign, giving a further positive mass after multiplication
by $\epsilon$. This proves the strict inequality.
The path and its primitive core class are supplied by the compact
spin realization of \cite[\S15]{WueUnramified}. Its ends are the equal
spin germ and a distinct-spin germ of the same positive sign. The
equal-spin germ has only the genus-$g$ elementary factor in degree
one. The displayed comparison concerns the endpoints; it requires
no assignment of a branched count to a single parameter of the path.
\end{proof}

One may express the result as a comparison of integral universal
products. Let $z$ be an indeterminate and let $Q_h(z)\in\ZZ[z]$ be the
polynomial with $Q_h(4\sin^2(t/2))=4\sin^2(ht/2)$. It is integral
by the cosine recurrence, with $Q_h(0)=0$ and $Q_1(z)=z$.
With $\mathcal P_w$ the partition set of \S\ref{p8-sec:characters} for the spin difference $w$, put
\[
 \mathcal Z(q,z)=\sum_{\lambda\in\mathcal P_w}
     (\epsilon q)^{|\lambda|}\prod_{s\in\lambda}Q_{h(s)}(z)^{g-1}.
\]
This has a unique formal expansion as an integral product of ordinary
balanced hook series of genera $h$ evaluated at $q^d$. Indeed, in
degree $d$ subtract the already determined factors of smaller degree
by multiplying by their inverses. The resulting coefficient is an
integral polynomial in $z$ divisible by $z^{d(g-1)}$. The degree-one
terms of the balanced factors are the distinct monomials $z^{h-1}$,
so its coefficients give the required integer exponents for this
degree. The divisibility is preserved under products and inversion,
and each degree requires only finitely many genera. Using these
exponents as in~\eqref{p8-eq:lifted-product} gives a balanced element
in the completed universal algebra.
\begin{corollary}\label{p8-cor:quotient}
This element agrees with $\widetilde P$ modulo the subgroup of
balanced products generated, with degree completion, by the
$F_h(q^d)$ with $h\ge H_d+2$.
\end{corollary}
\begin{proof}
Substitute $q=y t^{-2(g-1)}$ and $z=4\sin^2(t/2)$ in the formal
product. There are no negative powers of $t$ after this normalization.
The constant and quadratic coefficients in $t$ in its connected
logarithm are the unramified connected counts and
\eqref{p8-eq:character-first-excess}: a box of hook length $k$ contributes a
relative quadratic term $-(g-1)k^2/12$. These are exactly the
corresponding numerical coefficients of the geometric product.
Triangularity in covering degree and then in source genus, as in
\eqref{p8-eq:tails-exponent}, identifies its integer exponents in both
genera. Applying those exponents in the universal factors proves the
claim. Thus the two lowest genera can also be recovered directly from the collision and tail geometry, independently of the all-genus comparison in Theorem~\ref{thm:spin}.
\end{proof}

\section{Individual spin classes and their intersections}\label{sec:individual}
On covers with two simple branch points the residue description determines a class for each spin structure. Its character decomposition gives a lattice indexed by the two-torsion of the Jacobian. On covers with four simple branch points there is a further distinction: the Euler integral of two spins has contributions on several components of the stable-map stack. We isolate the component of finite smooth covers and compute the ten individual intersection pairings on its closure in the genus-two cubic case.
\subsection{The lattice of spin Euler classes}\label{p8-sec:lattice}
The individual Euler classes on the Hurwitz surface contain
information beyond their pairwise integrals. Fix a spin line
$\theta_0$, put $G=H^1(C;\ZZ/2)$, and use the flat sign line to
write $\theta_0 z$ for its twist by $z\in G$. On the complement of
a node's image use the permutation action of the normalization.
At a triple point its local monodromy is even. In either case the
permutation sign on the complement of the collision extends to a
class $\alpha\in G$. For a boundary divisor $B$ in
\eqref{p8-eq:intro-divisor} write this class as $\alpha_B$. It is
constant on $B$: locally the covering data away from a disk about
the moving point do not vary, and removing and filling the point
induces an isomorphism on the extended sign characters.
Define
\[
 D_\alpha=\sum_{B:\alpha_B=\alpha} s_{\theta_0}(B)[B],
 \quad
 \Lambda=\sum_{\alpha\in G}\ZZ D_\alpha,\qquad
 \Lambda_{\rm spin}=\sum_{z\in G} \ZZ\cdot4e(V_{\theta_0 z}),
\]
as lattices in $H^2(\mathcal F_m;\QQ)$. Here $V_\theta$ is the
oriented spin cokernel.
\begin{theorem}\label{p8-thm:lattice}
For every $m\ge2$, the $D_\alpha$ are independent, the spin Euler
classes span the same rational space, and
\[
     \Lambda/\Lambda_{\rm spin}
          \simeq \bigoplus_{j=1}^{2g}
               (\ZZ/2^j\ZZ)^{\binom{2g}{j}}.
\]
The least positive integer taking any one of the $e(V_\theta)$
into $\Lambda$ is four.
\end{theorem}
\begin{proof}
Every $\alpha$ occurs already on a nodal divisor. For nonzero
$\alpha$, choose surface generators with value one on $a_1$ alone.
If $m$ is even, send $a_1$ to an $m$-cycle. If $m$ is odd, send
$a_1$ to a transposition and $a_2$ to an $m$-cycle. The other chosen
generators are trivial, so this gives a connected unramified cover
with sign $\alpha$. For $\alpha=0$ and odd $m$ an $m$-cycle suffices.
For even $m\ge4$, assign an $(m-1)$-cycle and a $3$-cycle involving
one of its letters and the last sheet so that the generated action
is transitive; both permutations are even. In each instance glue a
pair of points in a fiber. For $m=2$ and $\alpha=0$ take two trivial
sheets and glue their points over the chosen value. Distinct reduced
boundary divisors are disjoint with strictly negative self-intersection.
The $D_\alpha$ therefore have mutually orthogonal nonzero squares.

Set $q_{\theta_0}(z)=h^0(\theta_0 z)-h^0(\theta_0)\bmod2$.
The ratio of the two boundary parities, by the orthogonal trace-bundle
calculation in Section~\ref{p8-sec:characters}, gives
\[
 -4e(V_{\theta_0 z})
    =(-1)^{m q_{\theta_0}(z)}
       \sum_{\alpha\in G}
          (-1)^{\langle z\smile\alpha,[C]\rangle}D_\alpha .
\]
The common row signs do not change the generated lattice. In dual
binary bases for the nondegenerate cup product the character matrix
is the $2g$-fold tensor power of the matrix with rows $(1,1)$ and
$(1,-1)$. Integral row and column operations diagonalize this
two-by-two matrix with entries $1,2$. Tensoring those unimodular
operations gives the Smith factors $2^j$ with the stated
multiplicities. Every coefficient of a single $e(V_\theta)$ in the
independent $D_\alpha$ is $\pm1/4$, proving the last assertion.
\end{proof}
The theorem compares two lattices generated by geometric classes in
rational cohomology. It does not require that the individual
boundary divisors span the integral cohomology of the stack.
\subsection{Four simple branch points}\label{sec:excess}
Fix two theta characteristics \(a,b\) on \(C\). Write \(c_d\) for the signed connected étale degree and \(F_d\) for the finite-map Euler integral on the first Hurwitz surface. Let \(M_d\) be the closure of the smooth connected covers with four distinct simple branch points, with the sum of its reduced four-dimensional component cycles. Write \(V_a,V_b\) for the real spin cokernels and put \(I_d=\int_{[M_d]}e(V_a)e(V_b)\).
\begin{proposition}\label{prop:excess}
For every \(d\ge2\), putting \(t=d(g-1)\), one has
\[
 L_{d,H_d+2}=I_d-\frac{t}{12}F_d+\frac{t(5t-1)}{1440}c_d.
\]
For a genus-two target and distinct even spins this gives \(I_2=0,\ L_{2,5}=-1/160,\ n_{2,5}=0\) and \(I_3=3/2\). For a repeated even spin on the same target, \(I_3=161/6\).
\end{proposition}
\begin{proof}
First remove the open and closed summands consisting of maps of degree one over a connected étale target cover of full degree \(d\); its lift is taken modulo its deck group. Their degree-one series contributes \([v^{2t+4}]x^t c_d=t(5t-1)c_d/1440\). Covers factoring through a smaller étale target remain. Normalize a map in the complement. The sum of half the ramification degree, the genera of contracted normalizations and the excess from the dual graph is two. The remaining positive-dimensional contracted strata have a pointed elliptic tail on a two-branch finite core, an elliptic component bridging two sheets, or a rational tripod. The first family has a proper parametrization by the universal pointed first Hurwitz surface times \(\overline{\mathcal M}_{1,1}\), including the inserted bridge when the mark meets a node. The degree on the generic image is one.

The other strata lie on \(M_d\): attaching a degree-two bubble by the pencil \(|p+q|\) smooths an elliptic bridge with attaching pair \(p,q\). For the tripod attach the rational function \(z+z^{-1}+(z-1)^{-1}\) at its three simple poles. It has four simple branch values, since the binary discriminant in the target coordinate is \(w^4-2w^3-5w^2+6w-23\), of discriminant \(-5619712\). At a direct node glue rational bubble covers of the matching contact orders and smooth the balanced equations after base change. In each case the stabilized target remains \(C\). The remaining support is the union of \(M_d\) and the elliptic-tail image, both pure of dimension four. A smooth simply branched cover has generic virtual multiplicity one, and at a generic tail the derivative at the unramified attaching point identifies the smoothing line with \(T_{f(p)}C\otimes H^1(E,\mathcal O_E)\). The generic virtual multiplicity there is also one. A cycle of the top dimension is determined by these multiplicities. On the tail the matching constants cancel, leaving \(\mathbb E^\vee\otimes\operatorname{ev}^*(a\oplus b)\) with top class \(-\lambda\operatorname{ev}^*c_1(K_C)\). Its integral along the two new parameters is \(-t/12\). This proves the formula; Calculation~B3 of \cite{WueSupplement} gives the full normalization along the lower-dimensional strata. Substitution of the degree-two and cubic coefficients of Theorem~\ref{thm:spin} and the signed finite sums of Sections~\ref{p8-sec:hurwitz}--\ref{p8-sec:characters} gives the stated values; for the repeated spin one uses the orientable balanced series. See also Calculations~B1--B4 of \cite{WueSupplement}.
\end{proof}

For genus two, distinct even spins give the entire connected double series
\[
 L_{2,k}=(-1)^k\frac{2^{2k-2}-4}{(2k-2)!},\qquad k\ge3 .
\]
Indeed the disconnected coefficient of \(Q^2\) vanishes and its connected coefficient is \(-x^2/2\). Subtracting the balanced factor in degree one leaves no elementary double exponent at genus five or higher. A separate geometric bound for reduced multisections is available. If the normal equations are independent line equations of degrees at most \(d_{\max}\), and the reduced graph is generically one-to-one without contracted components, the ramification degree plus twice the number of identified nodes is at most \(m(m-1)d_{\max}\). On the normalized ordered-pair cover the difference of the two solutions satisfies a real line equation with positive zero orders. Its zeros at ramification and at the ordered pairs identified by a node give the inequality by integrating the first Chern class. Pairs on which one component difference is identically zero are counted using the other line. In the triangular case one first separates distinct scalar projections and then applies the line bound in each equal-projection group. The projected-jet variation, the prime-degree consequence and the restriction at composite towers are detailed in Calculations~A1--A10 of \cite{WueSupplement}. In particular for noneffective spins and prime degree \(p\) the elementary genus is at most \(1+(g-1)\binom{p+1}{2}\).

\subsection{The ten even spins in genus two}
On \(M_3\subset\overline{\mathcal M}_6(C,3)\) put \(W_a=-R\pi_*f^*a\) and \(\Delta_a=(c_2(W_a)-e(V_a))/2\).
\begin{theorem}\label{thm:gram}
For every smooth genus-two target the ten even spins have \(\int_{[M_3]}\Delta_a^2=6\) and \(\int_{[M_3]}\Delta_a\Delta_b=-1/3\) for \(a\ne b\). The ten classes span a positive-definite space of rank ten.
\end{theorem}
\begin{proof}
Let \(\iota\) be the hyperelliptic involution of \(C\). For distinct even spins \(a,b\), the two residue-one Szeg\H o kernels have the same polar part along the diagonal. Their squared difference restricts there to a quadratic differential \(Q_{ab}\). The common section locus on the genus-five finite Hurwitz surface is its zero scheme. To see the condition, consider the trace-free part of \(f_*\mathcal O_Y\). The two pulled-back sections identify it with an elementary enlargement of \(a^{-1}\oplus b^{-1}\) at a point \(p\). In square frames at \(p\), the normalized meromorphic sections of the two spins have expansions \(z^{-1}+c_a z+\cdots\) and \(z^{-1}+c_b z+\cdots\). Expressing the cubic multiplication in the enlarged frame leaves one regularity condition, \(c_a=c_b\), or \(Q_{ab}(p)=0\). This calculation works over local Artin bases. The resulting cover is unique, smooth, with one triple ramification point over \(\iota(p)\); its monodromy is \(S_3\), so it has no deck factor. The full cubic multiplication and the boundary classification are in Calculation~C1 and Calculation~C2, \S\S2--3 of \cite{WueSupplement}. On the four-cycle \(M_3\) the common support consists of the disjoint curves \(E_p\simeq\overline{\mathcal M}_{1,1}\) obtained by attaching an elliptic tail at that triple point.

Suppose that \(p\) is a simple zero of \(Q_{ab}\). At the triple point the map deforms as \(w^3+tw+v\); let \(w=x\) mark the attachment and let \(\tau\) smooth the tail node. The versal equation is \(\tau B=0\), where \(B|_{\tau=0}=t+3x^2\) up to a unit, by the principal-part obstruction to smoothing. Thus \(B=0\) is a smooth chart of the main component and \(\tau=0\) cuts its tail divisor. On this divisor the ramification marking satisfies \(t=-3x^2\). Each supported correction \(\Delta_\theta\) is the class of its kernel surface with coefficient one, by complex-to-real determinant continuation in Calculation~C2, \S5 of \cite{WueSupplement}. On the finite surface the two kernel divisors meet with length one at this simple zero. Pullback under the degree-two ramification marking makes their intersection inside the tail divisor equal to \(2[E_p]\). Its normal line in the main component is the node-smoothing line, of first Chern class \(-\lambda\) on \(E_p\). The self-intersection formula for that divisor therefore gives
\[
 \operatorname{Contr}_{E_p}(\Delta_a\Delta_b)
     =-2\int_{\overline{\mathcal M}_{1,1}}\lambda=-\frac1{12}.
\]
The stack Hodge degree includes the pointed-tail automorphisms. Since \(K_C^2\) has degree four, four simple zeros give total mixed degree \(-1/3\).

Put \(S_a=\int\Delta_a^2\). The equal-spin integral is \(161/6\) from Proposition~\ref{prop:excess}; the distinct-spin integral is \(3/2\). The Chern class of the complex index is independent of \(a\) rationally, since spin differences are torsion. Subtracting one half the square of the difference of the two real Euler classes from one half the sum of their squares gives
\[
    \frac32=\frac{161}{6}-2(S_a+S_b)+4\int\Delta_a\Delta_b .
\]
It follows that \(S_a+S_b=12\) whenever \(Q_{ab}\) has four simple zeros.

Write the six hyperelliptic branch values as \(e_i\). Each spin is an unordered partition into two triples. Given two distinct partitions choose sides sharing two elements with monic root polynomial \(P\); the opposite sides share two with root polynomial \(R\). Call the exchanged values \(c,d\). The hyperelliptic Szeg\H o formula \cite[Appendix~C, (C.3)--(C.4)]{dhs} gives \(Q_{ab}=(c-d)(P'R-PR')(dx/y)^2/4\). The binary quadratic \(P'R-PR'\) has discriminant \(4\prod_{P(a_i)=R(b_j)=0}(a_i-b_j)\ne0\). Only \(c,d\) can be exceptional branch roots. At a fixed \(c\) this happens precisely when \(u_{a_1}+u_{a_2}=u_{b_1}+u_{b_2}\) for \(u_e=1/(c-e)\). Of five distinct \(u_e\) there are at most three such disjoint-pair equalities. Indeed project to distinct reals and let their successive gaps be \(A,B,C',D>0\). The five possible equations are \(A=C'\), \(A=C'+D\), \(A=D\), \(A+B=D\), \(B=D\); at most three coexist. The graph on ten spins with the four-simple-zero pairs as edges has at least \(45-18=27\) edges. At a fixed spin and branch value at most one omitted element realizes the equality, so every vertex has degree at least three. A disconnected graph with this minimum degree would have at most \(21\) edges, a bipartite graph at most \(25\). Connectedness and an odd cycle force \(S_a=6\) throughout. The displayed equation now gives \(-1/3\) for every distinct pair. The Gram matrix has eigenvalues \(3\) and \(19/3\), the latter with multiplicity nine.
\end{proof}

The local contribution at a zero of order \(j\) is \(-j/12\), including a nonsimple zero. Indeed the supported Euler products on the different closed elliptic-tail curves add by excision. At each simple zero the local calculation in the proof applies irrespective of the other zeros. One double zero must therefore contribute \(-1/6\) by the global mixed pairing and the two simple contributions. With two double zeros send their underlying branch values to \(0,\infty\). They are the ramification points of the map \(P/R\); its deck involution is \(z\mapsto-z\). Thus the common pairs are \(\{\pm u\}\), \(\{\pm v\}\). Writing \(k=uv\), the hyperelliptic curve has equation \(y^2=z(z^2-u^2)(z^2-v^2)\) and automorphism \((z,y)\mapsto(k/z,k^{3/2}y/z^3)\). It exchanges the two supports and preserves each spin because it sends its triple to its complement. Choose square-compatible lifts and invariant metrics by averaging over the finite extension by signs. The spin orientation continued from the complex determinant and the supported Euler product are then preserved. Hence the two contributions agree, each being \(-1/6\).

On \(y^2=\prod_{j=1}^6(x-j)\), the two exceptional pairs, indexing a spin by its triple containing \(1\), are \(126,135\) and \(135,156\). The smooth-source kernel pairing equals \(5\) for every even spin: the ramification-tail term equals \(1\) and the complete square is \(6\). Calculations~C3--C4 of \cite{WueSupplement} give the characteristic-degree calculation on the singular spectral incidence and its separate tail calculation. Identifying the two compactifications at cusps and tacnodes remains open.

\section{Triple ramification and multisections of a spin line}\label{p8-sec:graphs}

The triple-ramification divisors give a test for the appearance of
actual curves. A multisection here is the image of a curve in the
total space of a line or vector bundle whose projection to the base
is finite and surjective. We call its parametrization simple if it
has generic degree one onto its image.

\begin{theorem}\label{p8-thm:algebraic-graphs}
Let $C$ be any smooth projective curve of genus $g\ge2$, and let
$\theta^2=K_C$. On every connected component of the space of connected
degree-$m$ covers of $C$ with exactly one triple ramification point
and no other ramification, there is a cover $f:Y\to C$ such that
\[
                       h^0(Y,f^*\theta)>h^0(C,\theta).
\]
For every $m\ge3$ the holomorphic total space of $\theta$ consequently
contains simple multisections of degree $m$ with smooth normalization
of genus $H_m+1$ and a single triple ramification point over $C$.
Their supports can be chosen arbitrarily close to the zero section.
\end{theorem}
\begin{proof}
Let $T$ be such a connected component. It is a proper curve finite
in coarse space over $C$ by Section~\ref{p8-sec:hurwitz}. We may pass
to a connected finite level cover so that the family
$\pi:\mathcal Y\to T$ is a smooth projective family over a smooth
curve. For example, a symplectic level structure of level at least
three on the first homology of the smooth sources rigidifies their
automorphisms. Write $f:\mathcal Y\to C$ also for the evaluation map,
$p:T\to\mathcal Y$ for the ramification section, $E=p(T)$, and
$K=c_1(f^*K_C)$. The number $K\cdot E$ is strictly positive since
$T\to C$ moves the branch value. Relative Riemann--Hurwitz and the
first nonzero jet along $p$ give
\[
 c_1(\omega_\pi)=K+2E,\qquad K^2=0,\qquad E^2=-K\cdot E/3.
\]
Put $L=f^*\theta$. Grothendieck--Riemann--Roch for this smooth family
gives the degree of its negative determinant of cohomology as
\begin{align*}
 \deg c_1(-R\pi_*L)
 &=-\int_{\mathcal Y}
    \left(\tfrac12c_1(L)^2
       -\tfrac12 c_1(L)c_1(\omega_\pi)
       +\tfrac1{12}c_1(\omega_\pi)^2\right)\\
 &=\tfrac16 K\cdot E-\tfrac13 E^2=\tfrac5{18}K\cdot E .
\end{align*}

Suppose that every section on every fiber is a pullback. By cohomology
and base change $\pi_*L$ is then the constant bundle with fiber
$H^0(C,\theta)$. Relative duality and Riemann--Roch show that
$S=\pi_*(\omega_\pi L^{-1})/(df)\pi_* L$ is a line bundle, with
$\deg S^*=\deg c_1(-R\pi_*L)$. By the residue pairing $S$ is
isotropic in the rank-two orthogonal residue bundle. That bundle has
the global isotropic line $I=\pi_*(L(E)/L)$ of degree $K\cdot E/6$.
The two isotropic directions in an orthogonal plane are distinct; once
one is chosen, the other is a line paired perfectly with it.
Connectedness therefore forces $\deg S^*=\pm K\cdot E/6$.
Both values contradict the degree $5K\cdot E/18$. This argument
uses only the algebraic isotropics and includes effective $\theta$.

For each $m\ge3$ choose standard surface generators. On the first
handle use $(12),(23)$, and on the second use $(1\,2\,\ldots\,m)$
and the identity, with the remaining generators trivial. The
commutator product is a $3$-cycle, and these permutations generate
$S_m$. Filling the puncture gives a connected cover with the required
ramification. Its Hurwitz component consists of covers with the same
full monodromy. Take the additional section just proved to exist.
Its graph $(f,s)$ is simple: the normalization of its image gives an
intermediate cover, whereas the natural action of $S_m$ on $m$
sheets is primitive. A normalized image of degree one over $C$
would make $s$ a pullback, and the other possibility makes the graph
generically injective. The graphs of positive scalar multiples of
$s$ have different maximum fiber norms and converge to the zero
section. Their normalizations have the stated genus by
Riemann--Hurwitz.
\end{proof}

\begin{theorem}\label{p8-thm:path-graphs}
Let $a,b$ be distinct non-effective theta characteristics on $C$.
Identify their underlying smooth complex lines, and take any
continuous path of real Cauchy--Riemann operators on this line
between their spin operators. Suppose the operators on $C$ are
invertible throughout. In every degree $m\ge3$ there is an interior
parameter at which a smooth cover $f:Y\to C$ with full symmetric
monodromy and a single triple ramification point has nonzero
pullback kernel. Its kernel sections give simple almost-complex
multisections accumulating at the zero section. Such paths and
pairs exist in every base genus $g\ge2$. A single split path on
a closed symplectic six-torus realizes these multisections in all
degrees $m\ge3$.
\end{theorem}
\begin{proof}
Choose the first generator in the preceding permutation construction
to represent the Poincar\'e dual of the nonzero difference $ba^{-1}$
modulo two. Its monodromy is odd. On the corresponding compact
triple divisor $T$, the ratio of the two induced spin parities is
$-1$ by Section~\ref{p8-sec:characters}, since both embedded spin lines
have even parity. By Theorem~\ref{p8-thm:profiles} the oriented
obstruction planes at the two endpoints have opposite nonzero Euler
degrees on $T$. If every pullback along the path were injective
there, the constant-rank smooth-domain Fredholm family would give
an oriented plane bundle over $T\times[0,1]$. Its two Euler degrees
would be equal. This proves the existence of a nonzero section.
It cannot descend since the base operator has zero kernel. In using
primitivity for its graph, normalization of a simple image is the
usual factorization for a nonconstant pseudoholomorphic map.
The projection from the normalization to $C$ is holomorphic.
A degree-one projection is an isomorphism, and the resulting graph
over $C$ solves the base equation. This proves simplicity just as
in the holomorphic case. The positive real multiples again give
distinct supports. The assertion permits the same parameter to
serve more than one degree.

To construct a path, decrease the zeroth-order term of the spin
operator on $a$ to zero. Every strictly positive amplitude is
injective by the spin norm identity, and $\bar\partial_a$ is
invertible by non-effectivity and index zero. Join the complex
operators by a path in the complement of the theta divisor of
$\operatorname{Pic}^{g-1} C$, which is connected. A family of
representatives along the interval identifies the underlying
smooth lines. Reverse the amplitude deformation at $b$.
For existence in each genus take a hyperelliptic curve
$y^2=\prod_{i=1}^{2g+2}(x-e_i)$, with ramification points $W_i$
and divisor $H$ over infinity. For $|A|=g+1$ the line
$\mathcal O_C(\sum_{i\in A}W_i-H)$ squares to $K_C$ and has no
sections. Indeed, an invariant meromorphic function with at most
these simple poles and vanishing over infinity is zero. An
anti-invariant one is $yP(x)/\prod_{i\in A}(x-e_i)$ for a polynomial
$P$, and vanishing over infinity again forces zero. Replacing one
element of $A$ changes the line, since two distinct points on $C$
are not linearly equivalent. Finally keep one spin summand fixed
and put the path in the other summand of the normal bundle of
degree $2g-2$. The graph equation defines a path of almost-complex
germs with the prescribed projections. The compact realization
of \cite[Proposition~3.1 and \S15]{WueUnramified} extends the family from a fixed small
tube to tame structures on a symplectic six-torus. For each kernel
section a sufficiently small multiple lies in this tube, giving
a curve in the class $m[C]$ with the specified normalization.
\end{proof}

\begin{thebibliography}{Wue26e}
\bibitem[ACW17]{acw} D.~Abramovich, C.~Cadman and J.~Wise, \emph{Relative and orbifold Gromov--Witten invariants}, Algebr. Geom. \textbf{4} (2017), 472--500, \href{https://doi.org/10.14231/AG-2017-025}{doi:10.14231/AG-2017-025}.

\bibitem[AF16]{af} D.~Abramovich and B.~Fantechi, \emph{Orbifold techniques in degeneration formulas}, Ann. Sc. Norm. Super. Pisa Cl. Sci. (5) \textbf{16} (2016), 519--579, \href{https://doi.org/10.2422/2036-2145.201408_006}{doi:10.2422/2036-2145.201408\_006}.

\bibitem[AMW14]{amw} D.~Abramovich, S.~Marcus and J.~Wise, \emph{Comparison theorems for Gromov--Witten invariants of smooth pairs and of degenerations}, Ann. Inst. Fourier (Grenoble) \textbf{64} (2014), 1611--1667, \href{https://doi.org/10.5802/aif.2892}{doi:10.5802/aif.2892}.

\bibitem[Ati71]{ati} M.~F.~Atiyah,
\emph{Riemann surfaces and spin structures}, Ann. Sci. \'Ecole Norm. Sup. (4)
\textbf{4} (1971), 47--62, \href{https://doi.org/10.24033/asens.1205}{doi:10.24033/asens.1205}.

\bibitem[BS24]{bai} S.~Bai and M.~Swaminathan,
\emph{Bifurcations of embedded curves and toward an extension of Taubes's
Gromov invariant to Calabi--Yau $3$-folds}, Duke Math. J. \textbf{173}
(2024), 2947--3057,
\href{https://doi.org/10.1215/00127094-2023-0070}{doi:10.1215/00127094-2023-0070}.
Numbering as in \href{https://arxiv.org/abs/2106.01206v4}{arXiv:2106.01206v4}.

\bibitem[BS18]{baraglia} D.~Baraglia and L.~P.~Schaposnik,
\emph{Monodromy of rank $2$ twisted Hitchin systems and real character varieties},
Trans. Amer. Math. Soc. \textbf{370} (2018), 5491--5534,
\href{https://doi.org/10.1090/tran/7144}{doi:10.1090/tran/7144}.

\bibitem[BGG12]{bgg} S.~B.~Bradlow, O.~Garc\'ia-Prada and P.~B.~Gothen, \emph{Deformations of maximal representations in $\mathrm{Sp}(4,\mathbb R)$}, Q.~J.~Math. \textbf{63} (2012), 795--843, \href{https://arxiv.org/abs/0903.5496}{arXiv:0903.5496}.

\bibitem[BP01]{bp01} J.~Bryan and R.~Pandharipande,
\emph{BPS states of curves in Calabi--Yau $3$-folds},
Geom. Topol. \textbf{5} (2001), 287--318,
\href{https://doi.org/10.2140/gt.2001.5.287}{doi:10.2140/gt.2001.5.287}.

\bibitem[BP08]{bp08} J.~Bryan and R.~Pandharipande,
\emph{The local Gromov--Witten theory of curves}, J. Amer. Math. Soc.
\textbf{21} (2008), 101--136,
\href{https://doi.org/10.1090/S0894-0347-06-00545-5}{doi:10.1090/S0894-0347-06-00545-5}.

\bibitem[Cor89]{cornalba} M.~Cornalba, \emph{Moduli of curves and theta-characteristics},
in \emph{Lectures on Riemann surfaces} (Trieste, 1987), World Scientific, Teaneck, NJ, 1989, 560--589,
\href{https://mate.unipv.it/cornalba/papers/theta_moduli.pdf}{author's text}.

\bibitem[DS21]{dhs} E.~D'Hoker and O.~Schlotterer, \emph{Two-loop superstring five-point amplitudes. Part III. Construction via the RNS formulation: even spin structures}, J. High Energy Phys. \textbf{12} (2021), 063, \href{https://doi.org/10.1007/JHEP12(2021)063}{doi:10.1007/JHEP12(2021)063}.

\bibitem[DIW26]{diw} A.~Doan, E.-N.~Ionel and T.~Walpuski,
\emph{The Gopakumar--Vafa finiteness conjecture},
Ann. of Math. (2) \textbf{204} (2026), 819--871,
\href{https://doi.org/10.4007/annals.2026.204.3.1}{doi:10.4007/annals.2026.204.3.1}.
Numbering as in \href{https://arxiv.org/abs/2103.08221v2}{arXiv:2103.08221v2}.

\bibitem[EOP08]{eop} A.~Eskin, A.~Okounkov and R.~Pandharipande,
\emph{The theta characteristic of a branched covering}, Adv. Math.
\textbf{217} (2008), 873--888,
\href{https://doi.org/10.1016/j.aim.2006.08.001}{doi:10.1016/j.aim.2006.08.001}.

\bibitem[GI21]{gi} P.~Georgieva and E.-N.~Ionel,
\emph{A Klein TQFT: the local Real Gromov--Witten theory of curves},
Adv. Math. \textbf{391} (2021), Paper No.~107972,
\href{https://doi.org/10.1016/j.aim.2021.107972}{doi:10.1016/j.aim.2021.107972}.

\bibitem[IP18]{ip} E.-N.~Ionel and T.~H.~Parker,
\emph{The Gopakumar--Vafa formula for symplectic manifolds},
Ann. of Math. (2) \textbf{187} (2018), 1--64,
\href{https://arxiv.org/abs/1306.1516}{arXiv:1306.1516}.

\bibitem[Juc74]{jucys} A.-A.~A.~Jucys,
\emph{Symmetric polynomials and the center of the symmetric group ring},
Rep. Math. Phys. \textbf{5} (1974), 107--112,
\href{https://doi.org/10.1016/0034-4877(74)90019-6}{doi:10.1016/0034-4877(74)90019-6}.

\bibitem[LP12]{lpob} J.~Lee and T.~H.~Parker,
\emph{An obstruction bundle relating Gromov--Witten invariants of curves
and K\"ahler surfaces}, Amer. J. Math. \textbf{134} (2012), 453--506,
\href{https://doi.org/10.1353/ajm.2012.0010}{doi:10.1353/ajm.2012.0010}.

\bibitem[LP13]{lp} J.~Lee and T.~H.~Parker,
\emph{Spin Hurwitz numbers and the Gromov--Witten invariants of K\"ahler surfaces},
Comm. Anal. Geom. \textbf{21} (2013), 1015--1060,
\href{https://doi.org/10.4310/CAG.2013.v21.n5.a6}{doi:10.4310/CAG.2013.v21.n5.a6}.

\bibitem[Men25]{menon} K.~Menon,
\emph{Sum of squares of hook lengths and contents}, Australas. J. Combin.
\textbf{92} (2025), 70--79,
\href{https://ajc.maths.uq.edu.au/pdf/92/ajc_v92_p070.pdf}{journal text}.

\bibitem[Mum71]{mumford} D.~Mumford,
\emph{Theta characteristics of an algebraic curve}, Ann. Sci. \'Ecole
Norm. Sup. (4) \textbf{4} (1971), 181--192,
\href{https://doi.org/10.24033/asens.1209}{doi:10.24033/asens.1209}.

\bibitem[Par25]{pardon} J.~Pardon,
\emph{Universally counting curves in Calabi--Yau threefolds},
\href{https://arxiv.org/abs/2308.02948v3}{arXiv:2308.02948v3} (2025).

\bibitem[Wue26b]{WueProduct} B.~J.~Wuebben,
\emph{A geometric product formula for universal curve counting on Calabi--Yau threefolds}, 23 November 2026,
\href{https://doi.org/10.5281/zenodo.24455518}{doi:10.5281/zenodo.24455518}.

\bibitem[Wue26g]{WueSupplement} B.~J.~Wuebben,
\emph{Supplementary calculations for branched covers and spin structures in local curve counting}, Calculations A1--H3, 24 November 2026, archived with \emph{Branched triple covers and wall crossing for local curve counts},
\href{https://doi.org/10.5281/zenodo.24510974}{doi:10.5281/zenodo.24510974}.

\bibitem[Wue26j]{WueUnramified} B.~J.~Wuebben,
\emph{Unramified covers and wall crossing for local Gromov--Witten invariants}, 24 November 2026,
\href{https://doi.org/10.5281/zenodo.24510972}{doi:10.5281/zenodo.24510972}.

\bibitem[Zin11]{zinger} A.~Zinger,
\emph{A comparison theorem for Gromov--Witten invariants in the
symplectic category}, Adv. Math. \textbf{228} (2011), 535--574,
\href{https://doi.org/10.1016/j.aim.2011.05.021}{doi:10.1016/j.aim.2011.05.021}.
\end{thebibliography}

\end{document}
